\chapter{Concordances}\label{ch:concordances}

Les concordances renvoient des références des recueils et des bases aux numéros du catalogue. Elles sont produites à partir des champs de références de chaque notice. Les numéros des notices fusionnées (106, 282, 329) et jointes (394) restent cités, puisque leurs références renvoient au même monument que la notice de référence.

\section{Concordance générale}

{\small
% généré par scripts/04_generated.py
\begin{longtable}{@{}r r l l l@{}}\caption{Concordance générale : notice, année de classement, première référence ICUR, EDCS, EDB.}\label{tab:concgen}\\\toprule N° & Année & ICUR & EDCS & EDB\\\midrule\endfirsthead\toprule N° & Année & ICUR & EDCS & EDB\\\midrule\endhead\bottomrule\endfoot
\footnotesize 1 & 340 & ICUR I 1419 & EDCS-39000407 & EDB31183 \\
2 & 340 & ICUR III 8417 & EDCS-35802005 & EDB25692 \\
3 & 340 & ICUR IV 11089 & EDCS-37200371 & EDB39548 \\
4 & 341 & ICUR I 1420 & EDCS-34501143 & EDB31184 \\
5 & 342 & ICUR I 1251 & EDCS-39101961 & EDB30459 \\
6 & 342 & ICUR IV 9556 & EDCS-36101085 & EDB20872 \\
7 & 343 & ICUR I 264 & EDCS-38600270 & EDB27782 \\
8 & 343 & ICUR I 2080 & EDCS-25402137 & EDB15588 \\
9 & 343 & ICUR III 8720 & EDCS-35900965 & EDB18778 \\
10 & 344 & ICUR IV 11752 & EDCS-37400847 & EDB699 \\
11 & 345 & ICUR I 1104 & EDCS-38700644 & EDB30041 \\
12 & 352 & ICUR II 4798 & EDCS-36101279 & — \\
13 & 345 & ICUR II 4799 & EDCS-36101280 & — \\
14 & 345 & ICUR IV 11753 & EDCS-37400848 & EDB700 \\
15 & 346 & ICUR I 1930 & EDCS-20000191 & EDB27212 \\
16 & 346 & ICUR I 3126 & EDCS-39900588 & EDB33080 \\
17 & 346 & ICUR I 1421 & EDCS-39000408 & EDB31185 \\
18 & 346 & ICUR IV 11754 & EDCS-35700937 & EDB965 \\
19 & 347 & ICUR IV 10851 & EDCS-19000686 & EDB39161 \\
20 & 347 & ICUR X 26312 & EDCS-01000568 & EDB10655 \\
21 & 347 & ICUR I 4003 & EDCS-34801295 & EDB12186 \\
22 & 348 & ICUR I 887 & EDCS-24100635 & EDB29417 \\
23 & 348 & ICUR I 81 & EDCS-38600097 & EDB27454 \\
24 & 348 & ICUR I 318 & EDCS-38600315 & EDB27940 \\
25 & 348 & ICUR II 4800 & EDCS-36101281 & — \\
26 & 348 & ICUR II 5711 & EDCS-38302409 & — \\
27 & 348 & ICUR III 6498 & EDCS-34801298 & EDB20721 \\
28 & 348 & ICUR IV 11755 & EDCS-37400849 & EDB966 \\
29 & 348 & ICUR IV 11756 & EDCS-37400850 & EDB6728 \\
30 & 349 & ICUR III 8418 & EDCS-35802006 & EDB25693 \\
31 & 349 & ICUR III 8419 & EDCS-35802007 & EDB25694 \\
32 & 349 & ICUR IV 11090 & EDCS-37200372 & EDB39549 \\
33 & 349 & ICUR IV 12524 & EDCS-37801588 & EDB42440 \\
34 & 350 & ICUR I 2596 & EDCS-03900021 & EDB32543 \\
35 & 350 & ICUR I 1422 & EDCS-39000409 & EDB31186 \\
36 & 351 & ICUR I 19 & EDCS-18100517 & EDB27137 \\
37 & 352 & ICUR I 2081 & EDCS-39300305 & EDB32043 \\
38 & 352 & ICUR II 4241 & EDCS-33900148 & EDB17657 \\
39 & 352 & ICUR III 8140 & EDCS-35100216 & EDB24663 \\
40 & 353 & ICUR I 1252 & EDCS-39101962 & EDB30464 \\
41 & 355 & ICUR I 2769 & EDCS-39900318 & EDB32698 \\
42 & 355 & ICUR II 4268 & EDCS-33900175 & EDB8228 \\
43 & 355 & ICUR III 8420 & EDCS-31900329 & EDB25695 \\
44 & 355 & ICUR IV 11091 & EDCS-37200373 & EDB39550 \\
45 & 356 & ICUR II 4802 & EDCS-36101283 & — \\
46 & 356 & ICUR III 6499 & EDCS-34801299 & EDB19876 \\
47 & 357 & ICUR I 3166 & EDCS-40300313 & EDB33117 \\
48 & 357 & ICUR III 8722 & EDCS-35900967 & EDB20990 \\
49 & 358 & ICUR VII 17438 & EDCS-35200265 & EDB18146 \\
50 & 358 & ICUR I 1423 & EDCS-38801380 & EDB20348 \\
51 & 358 & ICUR I 1847 & EDCS-39200298 & EDB31763 \\
52 & 358 & ICUR II 6022 & EDCS-34200648 & EDB17462 \\
53 & 358 & ICUR IV 12420 & EDCS-37600944 & EDB41720 \\
54 & 359 & ICUR I 1424 & EDCS-39000410 & EDB31187 \\
55 & 359 & ICUR IV 11757 & EDCS-37400851 & EDB6729 \\
56 & 360 & ICUR II 5858 & EDCS-33100921 & EDB4428 \\
57 & 360 & ICUR III 8141 & EDCS-35100217 & EDB24664 \\
58 & 360 & ICUR IV 11093 & EDCS-37200375 & EDB39552 \\
59 & 361 & ICUR III 8142 & EDCS-35100218 & EDB24665 \\
60 & 361 & ICUR I 2979 & EDCS-39900423 & EDB32878 \\
61 & 361 & ICUR II 6023 & EDCS-34200649 & EDB17463 \\
62 & 362 & ICUR IV 12526 & EDCS-18600145 & EDB41423 \\
63 & 362 & ICUR I 82 & EDCS-38600098 & EDB27455 \\
64 & 362 & ICUR I 2805 & EDCS-11700571 & EDB32728 \\
65 & 362 & ICUR II 4804 & EDCS-36101285 & — \\
66 & 362 & ICUR III 6500 & EDCS-34801300 & EDB19877 \\
67 & 362 & ICUR III 8143 & EDCS-35100219 & EDB24666 \\
68 & 362 & ICUR III 8421 & EDCS-35802008 & EDB25696 \\
69 & 362 & ICUR IV 9557 & EDCS-36101086 & EDB20873 \\
70 & 362 & ICUR IV 9558 & EDCS-36101087 & EDB18727 \\
71 & 362 & ICUR IV 11094 & EDCS-37200376 & EDB39553 \\
72 & 362 & ICUR IV 11758 & EDCS-37400852 & EDB6730 \\
73 & 363 & ICUR I 1426 & EDCS-35200401 & EDB31189 \\
74 & 363 & ICUR I 966 & EDCS-38700473 & EDB29599 \\
75 & 363 & ICUR I 1425 & EDCS-39000411 & EDB31188 \\
76 & 363 & ICUR I 2082 & EDCS-39300306 & EDB32044 \\
77 & 363 & ICUR I 2083 & EDCS-39300307 & EDB32045 \\
78 & 363 & ICUR I 2806 & EDCS-11700572 & EDB32729 \\
79 & 363 & ICUR I 2807 & EDCS-11700573 & EDB32730 \\
80 & 363 & ICUR I 3167 & EDCS-40300314 & EDB33118 \\
81 & 363 & ICUR IV 9559 & EDCS-36101088 & EDB19258 \\
82 & 363 & ICUR IV 11759 & EDCS-37400853 & EDB2091 \\
83 & 364 & — & EDCS-59400075 & — \\
84 & 364 & ICUR I 2084 & EDCS-39300308 & EDB32048 \\
85 & 364 & ICUR I 2770 & EDCS-39900319 & EDB32699 \\
86 & 364 & ICUR II 5789 & EDCS-38302491 & EDB22178 \\
87 & 364 & ICUR II 6024 & EDCS-34200650 & EDB17464 \\
88 & 364 & ICUR III 7381 & EDCS-34900709 & EDB24129 \\
89 & 364 & ICUR IV 9560 & EDCS-36101089 & EDB19824 \\
90 & 364 & ICUR IV 11096 & EDCS-37200378 & EDB39555 \\
91 & 364 & ICUR IV 11760 & EDCS-37400854 & EDB2092 \\
92 & 365 & ICUR I 1189 & EDCS-38700730 & EDB30135 \\
93 & 365 & ICUR I 2085 & EDCS-39300309 & EDB32049 \\
94 & 365 & ICUR I 2695 & EDCS-39900259 & EDB32628 \\
95 & 365 & ICUR II 6025 & EDCS-34200651 & EDB17465 \\
96 & 365 & ICUR III 7382 & EDCS-34900710 & EDB24130 \\
97 & 365 & ICUR IV 9561 & EDCS-36101090 & EDB18365 \\
98 & — & — & EDCS-16300190 & EDB20076 \\
99 & 366 & ICUR III 8723 & EDCS-22300012 & EDB20991 \\
100 & 366 & ICUR I 1429 & EDCS-39000413 & EDB31193 \\
101 & 366 & ICUR II 4163 & EDCS-33900327 & EDB19222 \\
102 & 366 & ICUR II 4269 & EDCS-33900176 & EDB8229 \\
103 & 366 & ICUR IV 10853 & EDCS-36901720 & EDB39162 \\
104 & 366 & ICUR IV 11097 & EDCS-37200379 & EDB39556 \\
105 & 366 & ICUR IV 11763 & EDCS-35700942 & EDB1004 \\
106 & 367 & — & EDCS-16600352 & EDB13167 \\
107 & 367 & ICUR I 2792 & EDCS-39900335 & EDB32714 \\
108 & 367 & ICUR I 265 & EDCS-38600271 & EDB27783 \\
109 & 367 & ICUR I 896 & EDCS-38700412 & EDB29441 \\
110 & 367 & ICUR I 940 & EDCS-34700228 & EDB19681 \\
111 & 367 & ICUR I 1350 & EDCS-39000328 & EDB30676 \\
112 & 367 & ICUR I 1401 & EDCS-34501139 & EDB30801 \\
113 & 367 & ICUR I 1431 & EDCS-39000414 & EDB31195 \\
114 & 367 & ICUR I 2086 & EDCS-39300310 & EDB32051 \\
115 & 367 & ICUR I 3175 & EDCS-40300321 & EDB33127 \\
116 & 367 & ICUR II 4805 & EDCS-36101286 & — \\
117 & 367 & ICUR II 4807 & EDCS-36101288 & — \\
118 & 367 & ICUR II 6026 & EDCS-34200652 & EDB17466 \\
119 & 367 & ICUR II 6027 & EDCS-34200653 & EDB13166 \\
120 & 367 & ICUR II 6028 & EDCS-34200654 & EDB13167 \\
121 & 367 & ICUR III 6502 & EDCS-34801302 & EDB19879 \\
122 & 367 & ICUR III 8146 & EDCS-35700175 & EDB24669 \\
123 & 367 & ICUR III 9344b & EDCS-35900809 & EDB24921 \\
124 & 367 & ICUR IV 9563 & EDCS-36101092 & EDB18367 \\
125 & 368 & ICUR I 963 & EDCS-35200402 & EDB29591 \\
126 & 368 & ICUR I 479 & EDCS-38700887 & EDB28428 \\
127 & 368 & ICUR I 725 & EDCS-38701126 & EDB28952 \\
128 & 368 & ICUR I 2808 & EDCS-11700574 & EDB32731 \\
129 & 368 & ICUR II 4500 & EDCS-34100563 & EDB20916 \\
130 & 368 & ICUR II 5931 & EDCS-38100789 & EDB21154 \\
131 & 368 & ICUR III 8147 & EDCS-35700176 & EDB24670 \\
132 & 369 & — & EDCS-29900387 & — \\
133 & 369 & ICUR II 4165 & EDCS-33900077 & EDB19224 \\
134 & 369 & ICUR III 8148 & EDCS-35700177 & EDB24671 \\
135 & 369 & ICUR IV 11101 & EDCS-37200383 & EDB39560 \\
136 & 370 & ICUR I 3179 & EDCS-40300325 & EDB33131 \\
137 & 370 & ICUR I 442 & EDCS-38600420 & EDB28273 \\
138 & 370 & ICUR I 1351 & EDCS-39000329 & EDB30683 \\
139 & 370 & ICUR I 3180 & EDCS-40300326 & EDB33132 \\
140 & 370 & ICUR II 4501 & EDCS-34100564 & EDB20917 \\
141 & 370 & ICUR II 6030 & EDCS-34200656 & EDB9633 \\
142 & 370 & ICUR IV 12421 & EDCS-37600945 & EDB41696 \\
143 & 371 & ICUR I 1432 & EDCS-39000415 & EDB31196 \\
144 & 371 & ICUR I 1936 & EDCS-35300743 & EDB35584 \\
145 & 371 & ICUR I 2088 & EDCS-39300312 & EDB32053 \\
146 & 371 & ICUR I 2809 & EDCS-11700575 & EDB32732 \\
147 & 371 & ICUR I 3181 & EDCS-40300327 & EDB33133 \\
148 & 371 & ICUR II 4502 & EDCS-34100565 & EDB20918 \\
149 & 371 & ICUR II 5712 & EDCS-38302410 & — \\
150 & 371 & ICUR II 5933 & EDCS-38100790 & EDB21458 \\
151 & 371 & ICUR II 6031 & EDCS-34200657 & EDB17432 \\
152 & 371 & ICUR IV 11102 & EDCS-37200384 & EDB39561 \\
153 & 371 & ICUR IV 11103 & EDCS-37200385 & EDB39562 \\
154 & 371 & ICUR IV 11765 & EDCS-35700944 & EDB2099 \\
155 & 372 & ICUR I 83 & EDCS-38600099 & EDB27456 \\
156 & 372 & ICUR I 1433 & EDCS-39000416 & EDB31197 \\
157 & 372 & ICUR I 1434 & EDCS-39000417 & EDB31198 \\
158 & 372 & ICUR I 2090 & EDCS-39300313 & EDB32056 \\
159 & 372 & ICUR I 3183 & EDCS-40300329 & EDB33135 \\
160 & 372 & ICUR II 4748 & EDCS-34201004 & EDB13470 \\
161 & 372 & ICUR II 4808 & EDCS-36101289 & — \\
162 & 372 & ICUR III 6503 & EDCS-34801303 & EDB22061 \\
163 & 372 & ICUR IV 9565 & EDCS-36101094 & EDB18369 \\
164 & 372 & ICUR IV 11104 & EDCS-37200386 & EDB39563 \\
165 & 372 & ICUR IV 11613t & EDCS-37300523 & EDB40555 \\
166 & 373 & ICUR I 1254 & EDCS-39101964 & EDB30476 \\
167 & 373 & ICUR I 2091 & EDCS-39300314 & EDB32057 \\
168 & 373 & ICUR I 2092 & EDCS-34700234 & EDB11297 \\
169 & 373 & ICUR III 8724 & EDCS-35900968 & EDB20992 \\
170 & 373 & ICUR III 9344c & EDCS-35900810 & EDB24922 \\
171 & 373 & ICUR IV 11766 & EDCS-35700945 & EDB2100 \\
172 & 374 & ICUR I 1937 & EDCS-39200381 & EDB31855 \\
173 & 374 & ICUR I 1938 & EDCS-39200382 & EDB31856 \\
174 & 374 & ICUR IV 9566 & EDCS-35700797 & EDB18370 \\
175 & 374 & ICUR IV 9567 & EDCS-35700798 & EDB18371 \\
176 & 374 & ICUR IV 11105 & EDCS-37200387 & EDB39564 \\
177 & 374 & ICUR IV 11106 & EDCS-37200388 & EDB39565 \\
178 & 374 & ICUR IV 12243 & EDCS-37600678 & EDB8839 \\
179 & 375 & ICUR V 13109 & EDCS-19900430 & EDB5945 \\
180 & 375 & ICUR I 1847 & EDCS-39200296 & EDB31761 \\
181 & 375 & ICUR I 2093 & EDCS-39300315 & EDB32058 \\
182 & 375 & ICUR I 3184 & EDCS-40300330 & EDB33136 \\
183 & 375 & ICUR II 6009 & EDCS-37700196 & EDB21312 \\
184 & 375 & ICUR II 6032 & EDCS-34200658 & EDB17433 \\
185 & 375 & ICUR III 8725 & EDCS-35900969 & EDB20993 \\
186 & 375 & ICUR III 8726 & EDCS-35900970 & EDB20994 \\
187 & 375 & ICUR IV 11767 & EDCS-37400856 & EDB876 \\
188 & 375 & ICUR IV 12531 & EDCS-37801594 & EDB42433 \\
189 & 376 & ICUR VIII 20794 & EDCS-24100416 & EDB15528 \\
190 & 376 & ICUR I 3186 & EDCS-40300332 & EDB33138 \\
191 & 376 & ICUR IV 9568 & EDCS-36101095 & EDB18372 \\
192 & 376 & ICUR IV 9569 & EDCS-35700799 & EDB17958 \\
193 & 376 & ICUR IV 11769 & EDCS-37400858 & EDB878 \\
194 & 376 & ICUR IV 12244 & EDCS-37600679 & EDB5083 \\
195 & 377 & — & EDCS-09400838 & EDB33615 \\
196 & 377 & — & EDCS-02700284 & — \\
197 & 377 & ICUR I 84 & EDCS-38600100 & EDB27457 \\
198 & 377 & ICUR I 1435 & EDCS-39000418 & EDB31199 \\
199 & 377 & ICUR I 2094 & EDCS-39300316 & EDB32059 \\
200 & 377 & ICUR I 3188 & EDCS-40300334 & EDB33139 \\
201 & 377 & ICUR I 3189 & EDCS-40300335 & EDB33140 \\
202 & 377 & ICUR II 4815 & EDCS-36101297 & — \\
203 & 377 & ICUR II 5790 & EDCS-38302492 & EDB22177 \\
204 & 377 & ICUR III 8423 & EDCS-35802010 & EDB25698 \\
205 & 377 & ICUR IV 11107 & EDCS-37200389 & EDB39569 \\
206 & 370 & ICUR IV 11770 & EDCS-35700946 & EDB880 \\
207 & 378 & ICUR IV 9570 & EDCS-36101096 & EDB17959 \\
208 & 378 & ICUR IV 11771 & EDCS-37400859 & EDB1342 \\
209 & 378 & ICUR IV 11772 & EDCS-37400860 & EDB5268 \\
210 & 379 & ICUR I 3142 & EDCS-39900598 & EDB33094 \\
211 & 379 & ICUR I 3191 & EDCS-40300336 & EDB33143 \\
212 & 379 & ICUR II 4948 & EDCS-37700147 & — \\
213 & 380 & ICUR II 4166 & EDCS-19900369 & EDB19166 \\
214 & 380 & ICUR I 3194 & EDCS-35200283 & EDB33145 \\
215 & 380 & ICUR I 1438 & EDCS-39000421 & EDB31203 \\
216 & 380 & ICUR I 3192 & EDCS-38801427 & EDB17710 \\
217 & 380 & ICUR II 6033 & EDCS-34200659 & EDB17434 \\
218 & 380 & ICUR IV 11774 & EDCS-37400862 & EDB5272 \\
219 & 381 & ICUR I 1209 & EDCS-38700748 & EDB30156 \\
220 & 381 & ICUR I 1352 & EDCS-39000330 & EDB30687 \\
221 & 381 & ICUR I 2095 & EDCS-39300317 & EDB32075 \\
222 & 381 & ICUR I 3196 & EDCS-40300339 & EDB33147 \\
223 & 381 & ICUR II 6034 & EDCS-34200660 & EDB17435 \\
224 & 381 & ICUR IV 9571 & EDCS-36101097 & EDB17960 \\
225 & 381 & ICUR IV 9572 & EDCS-36101098 & EDB17961 \\
226 & 381 & ICUR IV 11109 & EDCS-37200391 & EDB39571 \\
227 & 381 & ICUR IV 12259e & EDCS-35700732 & EDB4654 \\
228 & 382 & ICUR II 6036 & EDCS-16600356 & EDB3284 \\
229 & 382 & ICUR I 2096 & EDCS-19400473 & EDB32081 \\
230 & 382 & ICUR VI 17248 & EDCS-32900465 & EDB12805 \\
231 & 382 & ICUR I 1439 & EDCS-39000422 & EDB31204 \\
232 & 382 & ICUR I 3127 & EDCS-39900589 & EDB33081 \\
233 & 382 & ICUR I 3197 & EDCS-40300340 & EDB33148 \\
234 & 382 & ICUR II 4167 & EDCS-33900078 & EDB16941 \\
235 & 382 & ICUR II 5791 & EDCS-38302493 & EDB22056 \\
236 & 382 & ICUR II 5995 & EDCS-38100823 & EDB20813 \\
237 & 382 & ICUR III 6504 & EDCS-34801304 & EDB22062 \\
238 & 382 & ICUR III 8150 & EDCS-35700179 & EDB24673 \\
239 & 382 & ICUR III 8151 & EDCS-35700180 & EDB24674 \\
240 & 382 & ICUR IV 12259a & EDCS-35700730 & EDB4650 \\
241 & 383 & — & EDCS-13700250 & EDB42969 \\
242 & 383 & ICUR I 504 & EDCS-38700907 & EDB28484 \\
243 & 383 & ICUR I 1169 & EDCS-38700710 & EDB30112 \\
244 & 383 & ICUR I 1440 & EDCS-34501146 & EDB31205 \\
245 & 383 & ICUR I 2097 & EDCS-39300318 & EDB32082 \\
246 & 383 & ICUR I 2098 & EDCS-30600273 & EDB35608 \\
247 & 383 & ICUR I 2944 & EDCS-39900394 & EDB32852 \\
248 & 383 & ICUR I 3043 & EDCS-39900484 & EDB32975 \\
249 & 383 & ICUR I 3198 & EDCS-33500350 & EDB35683 \\
250 & 383 & ICUR II 6038 & EDCS-16600357 & EDB8750 \\
251 & 383 & ICUR II 6039 & EDCS-34200665 & EDB8751 \\
252 & 383 & ICUR III 7383 & EDCS-34900711 & EDB24131 \\
253 & 383 & ICUR III 8152 & EDCS-35700181 & EDB24675 \\
254 & 383 & ICUR III 8425 & EDCS-35802012 & EDB25700 \\
255 & 383 & ICUR III 8728 & EDCS-35900972 & EDB20583 \\
256 & 383 & ICUR IV 11111 & EDCS-37200393 & EDB39573 \\
257 & 383 & ICUR IV 12245 & EDCS-37600680 & EDB5084 \\
258 & 383 & ICUR IV 12246 & EDCS-37600681 & EDB5345 \\
259 & 383 & ICUR IV 12534 & EDCS-37801595 & EDB42430 \\
260 & 384 & ICUR VII 19963 & EDCS-03900106 & EDB21504 \\
261 & 384 & ICUR I 3201 & EDCS-40300342 & EDB33151 \\
262 & 384 & ICUR II 4168 & EDCS-19900458 & EDB16942 \\
263 & 384 & ICUR I 62 & EDCS-38600081 & EDB27435 \\
264 & 384 & ICUR I 177 & EDCS-38600196 & EDB27580 \\
265 & 384 & ICUR I 331 & EDCS-38600325 & EDB27988 \\
266 & 384 & ICUR I 481 & EDCS-38700889 & EDB28430 \\
267 & 384 & ICUR I 1210 & EDCS-38700749 & EDB30157 \\
268 & 384 & ICUR I 2099 & EDCS-39300319 & EDB32083 \\
269 & 384 & ICUR I 3200 & EDCS-40300341 & EDB33150 \\
270 & 384 & ICUR II 4480 & EDCS-34100543 & EDB19138 \\
271 & 384 & ICUR II 4750 & EDCS-34201006 & EDB13472 \\
272 & 384 & ICUR II 6040 & EDCS-34200666 & EDB8752 \\
273 & 384 & ICUR III 8426 & EDCS-35802013 & EDB25701 \\
274 & 384 & ICUR IV 11112 & EDCS-37200394 & EDB39574 \\
275 & 385 & ICUR I 1441 & EDCS-34501147 & EDB31206 \\
276 & 385 & ICUR I 3202 & EDCS-40300343 & EDB33152 \\
277 & 385 & ICUR II 4481 & EDCS-34100544 & EDB19139 \\
278 & 383 & ICUR II 6037 & EDCS-34200663 & EDB3285 \\
279 & 385 & ICUR II 6041 & EDCS-34200667 & EDB8753 \\
280 & 385 & ICUR IV 9573 & EDCS-36101099 & EDB17962 \\
281 & 385 & ICUR IV 12247 & EDCS-37600682 & EDB5346 \\
282 & 386 & — & EDCS-16600245 & EDB11278 \\
283 & 386 & ICUR I 1257 & EDCS-39101967 & EDB30485 \\
284 & 386 & ICUR I 1353 & EDCS-39000331 & EDB30696 \\
285 & 386 & ICUR I 1443 & EDCS-39000424 & EDB31210 \\
286 & 386 & ICUR I 2101 & EDCS-30600274 & EDB35609 \\
287 & 386 & ICUR II 6042 & EDCS-34200668 & EDB11277 \\
288 & 386 & ICUR II 6043 & EDCS-34200669 & EDB11278 \\
289 & 386 & ICUR III 8154 & EDCS-35700183 & EDB24676 \\
290 & 386 & ICUR IV 11114 & EDCS-37200396 & EDB39576 \\
291 & 387 & ICUR I 1258 & EDCS-39101968 & EDB30488 \\
292 & 387 & ICUR I 1444 & EDCS-39000425 & EDB31208 \\
293 & 387 & ICUR I 3204 & EDCS-40300344 & EDB33154 \\
294 & 387 & ICUR III 8155 & EDCS-35700184 & EDB24677 \\
295 & 388 & ICUR VII 17494 & EDCS-03900101 & EDB21373 \\
296 & 388 & ICUR VII 17495 & EDCS-19301140 & EDB21374 \\
297 & 388 & ICUR I 3206 & EDCS-33500351 & EDB35685 \\
298 & 388 & ICUR II 5996 & EDCS-38100824 & EDB21298 \\
299 & 388 & ICUR II 6044 & EDCS-34200670 & EDB11279 \\
300 & 388 & ICUR II 6045 & EDCS-34200671 & EDB11280 \\
301 & 388 & ICUR II 6046 & EDCS-34200672 & EDB11281 \\
302 & 388 & ICUR IV 9575 & EDCS-36101100 & EDB17188 \\
303 & 388 & ICUR IV 12250 & EDCS-37600684 & EDB5349 \\
304 & 389 & — & EDCS-08900756 & EDB27361 \\
305 & 389 & — & EDCS-08300019 & EDB34442 \\
306 & 389 & ICUR V 13355 & EDCS-01000473 & EDB27045 \\
307 & 389 & — & EDCS-74700288 & EDB31230 \\
308 & 389 & ICUR I 2771 & EDCS-39900320 & EDB32700 \\
309 & 389 & ICUR I 3208 & EDCS-38100791 & EDB21459 \\
310 & 389 & ICUR II 4823 & EDCS-36101305 & EDB43052 \\
311 & 388 & ICUR III 6506 & EDCS-34801306 & EDB22064 \\
312 & 389 & ICUR III 6507 & EDCS-34801307 & EDB22065 \\
313 & 389 & ICUR IV 11117a & EDCS-37200399 & EDB39579 \\
314 & 390 & ICUR IV 11117b & EDCS-37200400 & EDB39579 \\
315 & 389 & ICUR IV 11118 & EDCS-37200401 & EDB39580 \\
316 & 389 & ICUR IV 11119 & EDCS-37200402 & EDB39581 \\
317 & 389 & ICUR IV 11120 & EDCS-37200403 & EDB39582 \\
318 & 389 & ICUR IV 12251 & EDCS-37600685 & EDB6536 \\
319 & 390 & ICUR I 1445 & EDCS-39000426 & EDB31209 \\
320 & 390 & ICUR I 2102 & EDCS-39300321 & EDB32090 \\
321 & 390 & ICUR II 4778c & EDCS-20100025 & EDB42643 \\
322 & 390 & ICUR II 6047 & EDCS-34200673 & EDB11282 \\
323 & 390 & ICUR II 6048 & EDCS-34200674 & EDB14639 \\
324 & 390 & ICUR II 6049 & EDCS-34200675 & EDB10125 \\
325 & 390 & ICUR III 8156 & EDCS-35700185 & EDB24678 \\
326 & 390 & ICUR III 9344a & EDCS-35900808 & EDB24920 \\
327 & 390 & ICUR IV 9576 & EDCS-36101101 & EDB17189 \\
328 & 381 & ICUR IV 9578 & EDCS-36101103 & EDB17191 \\
329 & 391 & — & EDCS-07400001 & EDB10126 \\
330 & 391 & ICUR VII 19974 & EDCS-19900403 & EDB20295 \\
331 & 391 & ICUR I 2103 & EDCS-39300322 & EDB32097 \\
332 & 391 & ICUR I 3209 & EDCS-40300347 & EDB33157 \\
333 & 391 & ICUR II 4169 & EDCS-33900079 & EDB16943 \\
334 & 391 & ICUR II 4751 & EDCS-34201007 & EDB13473 \\
335 & 391 & ICUR II 5935 & EDCS-38100792 & EDB21460 \\
336 & 391 & ICUR II 6050 & EDCS-34200676 & EDB10126 \\
337 & 391 & ICUR II 6051 & EDCS-34200677 & EDB14891 \\
338 & 391 & ICUR III 8157 & EDCS-35700186 & EDB24679 \\
339 & 391 & ICUR III 8158 & EDCS-35700187 & EDB24680 \\
340 & 391 & ICUR III 8428 & EDCS-35802015 & EDB25702 \\
341 & 391 & ICUR IV 11121 & EDCS-37200404 & EDB39583 \\
342 & 391 & ICUR IV 11122 & EDCS-37200405 & EDB39584 \\
343 & 392 & ICUR I 986 & EDCS-38700493 & EDB29619 \\
344 & 392 & ICUR I 1355 & EDCS-39000332 & EDB30697 \\
345 & 392 & ICUR I 2104 & EDCS-39300323 & EDB32098 \\
346 & 392 & ICUR I 2980 & EDCS-39900424 & EDB32879 \\
347 & 392 & ICUR II 4271a & EDCS-33900177 & EDB11967 \\
348 & 392 & ICUR II 4271b & EDCS-25401950 & EDB11272 \\
349 & 392 & ICUR II 6052 & EDCS-34200678 & EDB14892 \\
350 & 392 & ICUR III 6508 & EDCS-34801308 & EDB22066 \\
351 & 393 & — & EDCS-07000089 & EDB41240 \\
352 & 393 & ICUR I 727 & EDCS-35100467 & EDB28954 \\
353 & 393 & ICUR I 1448 & EDCS-30300391 & EDB35388 \\
354 & 393 & ICUR I 1940 & EDCS-39200384 & EDB31858 \\
355 & 393 & ICUR III 8159 & EDCS-35700188 & EDB24681 \\
356 & 393 & ICUR III 8160 & EDCS-35700189 & EDB24682 \\
357 & 393 & ICUR III 8429 & EDCS-43700747 & EDB25703 \\
358 & 393 & ICUR III 8431 & EDCS-35802018 & EDB24939 \\
359 & 393 & ICUR IV 11123 & EDCS-37200406 & EDB39585 \\
360 & 393 & ICUR IV 11124a & EDCS-37200407 & EDB746 \\
361 & 393 & ICUR IV 12422 & EDCS-37600946 & EDB41699 \\
362 & 394 & — & EDCS-13000396 & EDB43058 \\
363 & 394 & ICUR II 4487 & EDCS-34100550 & EDB31779 \\
364 & 394 & ICUR II 4503 & EDCS-34100566 & EDB20919 \\
365 & 394 & ICUR II 6460 & EDCS-34300812 & EDB20818 \\
366 & 394 & ICUR III 8648 & EDCS-35900847 & EDB19494 \\
367 & 394 & ICUR IV 11125 & EDCS-37200411 & EDB39589 \\
368 & 395 & — & EDCS-23702281 & EDB41941 \\
369 & 395 & — & EDCS-44100241 & EDB42064 \\
370 & 395 & ICUR I 1212 & EDCS-38700751 & EDB30159 \\
371 & 395 & ICUR I 2105 & EDCS-39300324 & EDB32099 \\
372 & 395 & ICUR I 3213 & EDCS-40300349 & EDB33161 \\
373 & 395 & ICUR II 4510 & EDCS-34100573 & EDB18940 \\
374 & 395 & ICUR II 4829 & EDCS-36101309 & — \\
375 & 395 & ICUR II 5861 & EDCS-38302567 & EDB22029 \\
376 & 395 & ICUR II 6053 & EDCS-34200679 & EDB14893 \\
377 & 395 & ICUR III 6509 & EDCS-34801309 & EDB22067 \\
378 & 395 & ICUR III 8161 & EDCS-35700190 & EDB25348 \\
379 & 395 & ICUR III 8162 & EDCS-35700191 & EDB24684 \\
380 & 395 & ICUR III 8163 & EDCS-35700192 & EDB24685 \\
381 & 395 & ICUR III 8164 & EDCS-35700193 & EDB24686 \\
382 & 395 & ICUR III 8729 & EDCS-31900337 & EDB20584 \\
383 & 395 & ICUR IV 9377 & EDCS-36100891 & EDB19574 \\
384 & 395 & ICUR IV 11127 & EDCS-37200414 & EDB39591 \\
385 & 396 & — & EDCS-75200321 & — \\
386 & 396 & ICUR VII 17510 & EDCS-03900102 & EDB17954 \\
387 & 396 & — & EDCS-71500011 & — \\
388 & 396 & ICUR I 1450 & EDCS-39000429 & EDB31239 \\
389 & 396 & ICUR I 2106 & EDCS-39300325 & EDB32100 \\
390 & 396 & ICUR I 2107 & EDCS-39300326 & EDB32101 \\
391 & 396 & ICUR I 4073 & EDCS-40301037 & EDB34240 \\
392 & 396 & ICUR III 6510 & EDCS-34801310 & EDB22305 \\
393 & 396 & ICUR IV 11126a & EDCS-37200412 & EDB39590 \\
394 & 396 & ICUR IV 11126b & EDCS-37200413 & EDB39590 \\
395 & 396 & ICUR IV 12252a & EDCS-35700729 & EDB7067 \\
396 & 397 & — & EDCS-02700238 & EDB42676 \\
397 & 397 & ICUR I 3214 & EDCS-23102819 & EDB35690 \\
398 & 397 & ICUR I 471 & EDCS-38700879 & EDB28397 \\
399 & 397 & ICUR I 717 & EDCS-38701118 & EDB28944 \\
400 & 397 & ICUR I 941 & EDCS-38700455 & EDB29526 \\
401 & 397 & ICUR I 1451 & EDCS-39000430 & EDB31240 \\
402 & 397 & ICUR I 2108 & EDCS-39300327 & EDB32105 \\
403 & 397 & ICUR I 2109 & EDCS-39300328 & EDB32106 \\
404 & 397 & ICUR I 2772 & EDCS-39900321 & EDB32701 \\
405 & 397 & ICUR I 2812 & EDCS-11700578 & EDB32736 \\
406 & 397 & ICUR II 4830 & EDCS-36101310 & — \\
407 & 397 & ICUR II 4832 & EDCS-36101312 & — \\
408 & 397 & ICUR II 6055 & EDCS-34200681 & EDB16486 \\
409 & 397 & ICUR II 6056 & EDCS-34200682 & EDB16487 \\
410 & 397 & ICUR II 6057 & EDCS-34200683 & EDB16488 \\
411 & 397 & ICUR II 6058 & EDCS-34200684 & EDB16489 \\
412 & 397 & ICUR IV 9581 & EDCS-36101106 & EDB20443 \\
413 & 397 & ICUR IV 9582 & EDCS-36101107 & EDB19321 \\
414 & 397 & ICUR IV 11128 & EDCS-37200415 & EDB39592 \\
415 & 397 & ICUR IV 11778 & EDCS-37400866 & EDB5167 \\
416 & 397 & ICUR IV 11779 & EDCS-37400867 & EDB5168 \\
417 & 398 & — & EDCS-08400008 & EDB38747 \\
418 & 398 & ICUR I 1452 & EDCS-39000431 & EDB31242 \\
419 & 398 & ICUR I 1941 & EDCS-39200385 & EDB31859 \\
420 & 398 & ICUR II 4505 & EDCS-34100568 & EDB19792 \\
421 & 398 & ICUR III 8165 & EDCS-35700194 & EDB24687 \\
422 & 398 & ICUR III 8173b & EDCS-35700212 & EDB24704 \\
423 & 398 & ICUR III 8432 & EDCS-35802019 & EDB24940 \\
424 & 398 & ICUR III 8647 & EDCS-35900846 & EDB19129 \\
425 & 398 & ICUR IV 9583 & EDCS-36101108 & EDB19322 \\
426 & 398 & ICUR IV 11129 & EDCS-37200416 & EDB39593 \\
427 & 398 & ICUR IV 12424 & EDCS-37600947 & EDB41702 \\
428 & 398 & ICUR IV 12543 & EDCS-37801603 & EDB41825 \\
429 & 399 & — & EDCS-18200424 & — \\
430 & 399 & ICUR I 987 & EDCS-19900450 & EDB29620 \\
431 & 399 & ICUR I 309 & EDCS-38600308 & EDB27849 \\
432 & 399 & ICUR I 524 & EDCS-38700928 & EDB28503 \\
433 & 399 & ICUR I 1453 & EDCS-35100468 & EDB31243 \\
434 & 397 & ICUR II 6476 & EDCS-34300830 & EDB34468 \\
435 & 399 & ICUR III 6511 & EDCS-34801311 & EDB22306 \\
436 & 399 & ICUR III 8166 & EDCS-35700195 & EDB24688 \\
437 & 399 & ICUR IV 11124d & EDCS-37200410 & EDB39588 \\
438 & 399 & ICUR IV 11130 & EDCS-37200417 & EDB39594 \\
439 & 400 & ICUR I 3215 & EDCS-05801896 & EDB33162 \\
440 & 400 & ICUR I 1454 & EDCS-35200404 & EDB31244 \\
441 & 400 & ICUR I 528 & EDCS-38700931 & EDB28506 \\
442 & 400 & ICUR I 2110 & EDCS-39300329 & EDB32110 \\
443 & 400 & ICUR II 6059 & EDCS-34200685 & EDB16490 \\
444 & 400 & ICUR II 6060 & EDCS-34200686 & EDB16491 \\
445 & 400 & ICUR III 8168a & EDCS-35700197 & EDB24690 \\
446 & 400 & ICUR III 8730 & EDCS-35900973 & EDB20586 \\
447 & 401 & ICUR I 1942 & EDCS-25402147 & EDB12500 \\
448 & 401 & — & EDCS-23702282 & EDB41942 \\
449 & 401 & ICUR III 9365 & EDCS-35900834 & EDB20435 \\
450 & 401 & ICUR I 349 & EDCS-38600338 & EDB28045 \\
451 & 401 & ICUR II 4506 & EDCS-34100569 & EDB19793 \\
452 & 401 & ICUR II 4836 & EDCS-36101314 & — \\
453 & 401 & ICUR II 6061 & EDCS-34200687 & EDB13893 \\
454 & 401 & ICUR II 6062 & EDCS-34200688 & EDB13894 \\
455 & 401 & ICUR II 6063 & EDCS-34200689 & EDB13895 \\
456 & 401 & ICUR III 8167 & EDCS-35700196 & EDB24689 \\
457 & 401 & ICUR III 8168b & EDCS-35700198 & EDB24691 \\
458 & 401 & ICUR IV 11133 & EDCS-37200420 & EDB39597 \\
459 & 401 & ICUR IV 11134 & EDCS-37200421 & EDB39598 \\
460 & 402 & ICUR I 718 & EDCS-38701119 & EDB28945 \\
461 & 402 & ICUR I 1214 & EDCS-38700752 & EDB30165 \\
462 & 402 & ICUR I 1256 & EDCS-39101966 & EDB35372 \\
463 & 402 & ICUR I 1943 & EDCS-39200386 & EDB31861 \\
464 & 402 & ICUR I 3219 & EDCS-40300353 & EDB33166 \\
465 & 402 & ICUR I 4054 & EDCS-40301026 & EDB34225 \\
466 & 402 & ICUR II 6064 & EDCS-34200690 & EDB13896 \\
467 & 402 & ICUR III 8433 & EDCS-35802020 & EDB24941 \\
468 & 402 & ICUR IV 12253 & EDCS-37600687 & EDB7068 \\
469 & 403 & ICUR VIII 23434 & EDCS-23201556 & EDB21579 \\
470 & 403 & ICUR I 713 & EDCS-35100466 & EDB28939 \\
471 & 403 & ICUR I 1459 & EDCS-39000433 & EDB31246 \\
472 & 403 & ICUR I 1460 & EDCS-39000434 & EDB31247 \\
473 & 403 & ICUR I 3220 & EDCS-40300354 & EDB33167 \\
474 & 403 & ICUR I 3221 & EDCS-33900149 & EDB35329 \\
475 & 403 & ICUR II 4272 & EDCS-33900178 & EDB11273 \\
476 & 403 & ICUR II 4507 & EDCS-34100570 & EDB19794 \\
477 & 403 & ICUR IV 11136 & EDCS-37200423 & EDB39600 \\
478 & 403 & ICUR IV 12425 & EDCS-37600948 & EDB41721 \\
\end{longtable}

}

\section{Inscriptiones christianae urbis Romae}

% généré par scripts/04_generated.py
\begin{multicols}{4}\footnotesize\noindent ICUR I 19 → 36\\
ICUR I 62 → 263\\
ICUR I 81 → 23\\
ICUR I 82 → 63\\
ICUR I 83 → 155\\
ICUR I 84 → 197\\
ICUR I 177 → 264\\
ICUR I 264 → 7\\
ICUR I 265 → 108\\
ICUR I 309 → 431\\
ICUR I 318 → 24\\
ICUR I 331 → 265\\
ICUR I 349 → 450\\
ICUR I 442 → 137\\
ICUR I 471 → 398\\
ICUR I 479 → 126\\
ICUR I 481 → 266\\
ICUR I 504 → 242\\
ICUR I 524 → 432\\
ICUR I 528 → 441\\
ICUR I 713 → 470\\
ICUR I 717 → 399\\
ICUR I 718 → 460\\
ICUR I 725 → 127\\
ICUR I 727 → 352\\
ICUR I 887 → 22\\
ICUR I 896 → 109\\
ICUR I 940 → 110\\
ICUR I 941 → 400\\
ICUR I 963 → 125\\
ICUR I 966 → 74\\
ICUR I 986 → 343\\
ICUR I 987 → 430\\
ICUR I 1104 → 11\\
ICUR I 1169 → 243\\
ICUR I 1189 → 92\\
ICUR I 1209 → 219\\
ICUR I 1210 → 267\\
ICUR I 1212 → 370\\
ICUR I 1214 → 461\\
ICUR I 1251 → 5\\
ICUR I 1252 → 40\\
ICUR I 1254 → 166\\
ICUR I 1256 → 462\\
ICUR I 1257 → 283\\
ICUR I 1258 → 291\\
ICUR I 1350 → 111\\
ICUR I 1351 → 138\\
ICUR I 1352 → 220\\
ICUR I 1353 → 284\\
ICUR I 1355 → 344\\
ICUR I 1401 → 112\\
ICUR I 1419 → 1\\
ICUR I 1420 → 4\\
ICUR I 1421 → 17\\
ICUR I 1422 → 35\\
ICUR I 1423 → 50\\
ICUR I 1424 → 54\\
ICUR I 1425 → 75\\
ICUR I 1426 → 73\\
ICUR I 1429 → 100\\
ICUR I 1431 → 113\\
ICUR I 1432 → 143\\
ICUR I 1433 → 156\\
ICUR I 1434 → 157\\
ICUR I 1435 → 198\\
ICUR I 1438 → 215\\
ICUR I 1439 → 231\\
ICUR I 1440 → 244\\
ICUR I 1441 → 275\\
ICUR I 1443 → 285\\
ICUR I 1444 → 292\\
ICUR I 1445 → 319\\
ICUR I 1448 → 353\\
ICUR I 1450 → 388\\
ICUR I 1451 → 401\\
ICUR I 1452 → 418\\
ICUR I 1453 → 433\\
ICUR I 1454 → 440\\
ICUR I 1459 → 471\\
ICUR I 1460 → 472\\
ICUR I 1847 → 51\\
ICUR I 1847 → 180\\
ICUR I 1930 → 15\\
ICUR I 1936 → 144\\
ICUR I 1937 → 172\\
ICUR I 1938 → 173\\
ICUR I 1940 → 354\\
ICUR I 1941 → 419\\
ICUR I 1942 → 447\\
ICUR I 1943 → 463\\
ICUR I 2080 → 8\\
ICUR I 2081 → 37\\
ICUR I 2082 → 76\\
ICUR I 2083 → 77\\
ICUR I 2084 → 84\\
ICUR I 2085 → 93\\
ICUR I 2086 → 114\\
ICUR I 2088 → 145\\
ICUR I 2090 → 158\\
ICUR I 2091 → 167\\
ICUR I 2092 → 168\\
ICUR I 2093 → 181\\
ICUR I 2094 → 199\\
ICUR I 2095 → 221\\
ICUR I 2096 → 229\\
ICUR I 2097 → 245\\
ICUR I 2098 → 246\\
ICUR I 2099 → 268\\
ICUR I 2101 → 286\\
ICUR I 2102 → 320\\
ICUR I 2103 → 331\\
ICUR I 2104 → 345\\
ICUR I 2105 → 371\\
ICUR I 2106 → 389\\
ICUR I 2107 → 390\\
ICUR I 2108 → 402\\
ICUR I 2109 → 403\\
ICUR I 2110 → 442\\
ICUR I 2596 → 34\\
ICUR I 2695 → 94\\
ICUR I 2769 → 41\\
ICUR I 2770 → 85\\
ICUR I 2771 → 308\\
ICUR I 2772 → 404\\
ICUR I 2792 → 107\\
ICUR I 2805 → 64\\
ICUR I 2806 → 78\\
ICUR I 2807 → 79\\
ICUR I 2808 → 128\\
ICUR I 2809 → 146\\
ICUR I 2812 → 405\\
ICUR I 2944 → 247\\
ICUR I 2979 → 60\\
ICUR I 2980 → 346\\
ICUR I 3043 → 248\\
ICUR I 3126 → 16\\
ICUR I 3127 → 232\\
ICUR I 3142 → 210\\
ICUR I 3166 → 47\\
ICUR I 3167 → 80\\
ICUR I 3175 → 115\\
ICUR I 3179 → 136\\
ICUR I 3180 → 139\\
ICUR I 3181 → 147\\
ICUR I 3183 → 159\\
ICUR I 3184 → 182\\
ICUR I 3186 → 190\\
ICUR I 3188 → 200\\
ICUR I 3189 → 201\\
ICUR I 3191 → 211\\
ICUR I 3192 → 216\\
ICUR I 3194 → 214\\
ICUR I 3196 → 222\\
ICUR I 3197 → 233\\
ICUR I 3198 → 249\\
ICUR I 3200 → 269\\
ICUR I 3201 → 261\\
ICUR I 3202 → 276\\
ICUR I 3204 → 293\\
ICUR I 3206 → 297\\
ICUR I 3208 → 309\\
ICUR I 3209 → 332\\
ICUR I 3213 → 372\\
ICUR I 3214 → 397\\
ICUR I 3215 → 439\\
ICUR I 3219 → 464\\
ICUR I 3220 → 473\\
ICUR I 3221 → 474\\
ICUR I 4003 → 21\\
ICUR I 4054 → 465\\
ICUR I 4073 → 391\\
ICUR II 4163 → 101\\
ICUR II 4165 → 133\\
ICUR II 4166 → 213\\
ICUR II 4167 → 234\\
ICUR II 4168 → 262\\
ICUR II 4169 → 333\\
ICUR II 4241 → 38\\
ICUR II 4242 → 474\\
ICUR II 4268 → 42\\
ICUR II 4269 → 102\\
ICUR II 4271a → 347\\
ICUR II 4271b → 348\\
ICUR II 4272 → 475\\
ICUR II 4480 → 270\\
ICUR II 4481 → 277\\
ICUR II 4487 → 363\\
ICUR II 4500 → 129\\
ICUR II 4501 → 140\\
ICUR II 4502 → 148\\
ICUR II 4503 → 364\\
ICUR II 4505 → 420\\
ICUR II 4506 → 451\\
ICUR II 4507 → 476\\
ICUR II 4510 → 373\\
ICUR II 4748 → 160\\
ICUR II 4750 → 271\\
ICUR II 4751 → 334\\
ICUR II 4778c → 321\\
ICUR II 4798 → 12\\
ICUR II 4799 → 13\\
ICUR II 4800 → 25\\
ICUR II 4802 → 45\\
ICUR II 4804 → 65\\
ICUR II 4805 → 116\\
ICUR II 4807 → 117\\
ICUR II 4808 → 161\\
ICUR II 4815 → 202\\
ICUR II 4823 → 310\\
ICUR II 4829 → 374\\
ICUR II 4830 → 406\\
ICUR II 4832 → 407\\
ICUR II 4836 → 452\\
ICUR II 4948 → 212\\
ICUR II 5711 → 26\\
ICUR II 5712 → 149\\
ICUR II 5789 → 86\\
ICUR II 5790 → 203\\
ICUR II 5791 → 235\\
ICUR II 5858 → 56\\
ICUR II 5861 → 375\\
ICUR II 5931 → 130\\
ICUR II 5932 → 56\\
ICUR II 5933 → 150\\
ICUR II 5934 → 309\\
ICUR II 5935 → 335\\
ICUR II 5995 → 236\\
ICUR II 5996 → 298\\
ICUR II 6009 → 183\\
ICUR II 6022 → 52\\
ICUR II 6023 → 61\\
ICUR II 6024 → 87\\
ICUR II 6025 → 95\\
ICUR II 6026 → 118\\
ICUR II 6027 → 119\\
ICUR II 6028 → 120\\
ICUR II 6030 → 141\\
ICUR II 6031 → 151\\
ICUR II 6032 → 184\\
ICUR II 6033 → 217\\
ICUR II 6034 → 223\\
ICUR II 6036 → 228\\
ICUR II 6037 → 278\\
ICUR II 6038 → 250\\
ICUR II 6039 → 251\\
ICUR II 6040 → 272\\
ICUR II 6041 → 279\\
ICUR II 6042 → 287\\
ICUR II 6043 → 288\\
ICUR II 6044 → 299\\
ICUR II 6045 → 300\\
ICUR II 6046 → 301\\
ICUR II 6047 → 322\\
ICUR II 6048 → 323\\
ICUR II 6049 → 324\\
ICUR II 6050 → 336\\
ICUR II 6051 → 337\\
ICUR II 6052 → 349\\
ICUR II 6053 → 376\\
ICUR II 6055 → 408\\
ICUR II 6056 → 409\\
ICUR II 6057 → 410\\
ICUR II 6058 → 411\\
ICUR II 6059 → 443\\
ICUR II 6060 → 444\\
ICUR II 6061 → 453\\
ICUR II 6062 → 454\\
ICUR II 6063 → 455\\
ICUR II 6064 → 466\\
ICUR II 6460 → 365\\
ICUR II 6476 → 434\\
ICUR III 6498 → 27\\
ICUR III 6499 → 46\\
ICUR III 6500 → 66\\
ICUR III 6502 → 121\\
ICUR III 6503 → 162\\
ICUR III 6504 → 237\\
ICUR III 6506 → 311\\
ICUR III 6507 → 312\\
ICUR III 6508 → 350\\
ICUR III 6509 → 377\\
ICUR III 6510 → 392\\
ICUR III 6511 → 435\\
ICUR III 7381 → 88\\
ICUR III 7382 → 96\\
ICUR III 7383 → 252\\
ICUR III 8140 → 39\\
ICUR III 8141 → 57\\
ICUR III 8142 → 59\\
ICUR III 8143 → 67\\
ICUR III 8146 → 122\\
ICUR III 8147 → 131\\
ICUR III 8148 → 134\\
ICUR III 8150 → 238\\
ICUR III 8151 → 239\\
ICUR III 8152 → 253\\
ICUR III 8154 → 289\\
ICUR III 8155 → 294\\
ICUR III 8156 → 325\\
ICUR III 8157 → 338\\
ICUR III 8158 → 339\\
ICUR III 8159 → 355\\
ICUR III 8160 → 356\\
ICUR III 8161 → 378\\
ICUR III 8162 → 379\\
ICUR III 8163 → 380\\
ICUR III 8164 → 381\\
ICUR III 8165 → 421\\
ICUR III 8166 → 436\\
ICUR III 8167 → 456\\
ICUR III 8168a → 445\\
ICUR III 8168b → 457\\
ICUR III 8173b → 422\\
ICUR III 8345 → 378\\
ICUR III 8417 → 2\\
ICUR III 8418 → 30\\
ICUR III 8419 → 31\\
ICUR III 8420 → 43\\
ICUR III 8421 → 68\\
ICUR III 8423 → 204\\
ICUR III 8425 → 254\\
ICUR III 8426 → 273\\
ICUR III 8428 → 340\\
ICUR III 8429 → 357\\
ICUR III 8431 → 358\\
ICUR III 8432 → 423\\
ICUR III 8433 → 467\\
ICUR III 8647 → 424\\
ICUR III 8648 → 366\\
ICUR III 8720 → 9\\
ICUR III 8722 → 48\\
ICUR III 8723 → 99\\
ICUR III 8724 → 169\\
ICUR III 8725 → 185\\
ICUR III 8726 → 186\\
ICUR III 8728 → 255\\
ICUR III 8729 → 382\\
ICUR III 8730 → 446\\
ICUR III 9344a → 326\\
ICUR III 9344b → 123\\
ICUR III 9344c → 170\\
ICUR III 9365 → 449\\
ICUR IV 9377 → 383\\
ICUR IV 9556 → 6\\
ICUR IV 9557 → 69\\
ICUR IV 9558 → 70\\
ICUR IV 9559 → 81\\
ICUR IV 9560 → 89\\
ICUR IV 9561 → 97\\
ICUR IV 9563 → 124\\
ICUR IV 9565 → 163\\
ICUR IV 9566 → 174\\
ICUR IV 9567 → 175\\
ICUR IV 9568 → 191\\
ICUR IV 9569 → 192\\
ICUR IV 9570 → 207\\
ICUR IV 9571 → 224\\
ICUR IV 9572 → 225\\
ICUR IV 9573 → 280\\
ICUR IV 9575 → 302\\
ICUR IV 9576 → 327\\
ICUR IV 9578 → 328\\
ICUR IV 9581 → 412\\
ICUR IV 9582 → 413\\
ICUR IV 9583 → 425\\
ICUR IV 10851 → 19\\
ICUR IV 10853 → 103\\
ICUR IV 11089 → 3\\
ICUR IV 11090 → 32\\
ICUR IV 11091 → 44\\
ICUR IV 11093 → 58\\
ICUR IV 11094 → 71\\
ICUR IV 11096 → 90\\
ICUR IV 11097 → 104\\
ICUR IV 11101 → 135\\
ICUR IV 11102 → 152\\
ICUR IV 11103 → 153\\
ICUR IV 11104 → 164\\
ICUR IV 11105 → 176\\
ICUR IV 11106 → 177\\
ICUR IV 11107 → 205\\
ICUR IV 11109 → 226\\
ICUR IV 11111 → 256\\
ICUR IV 11112 → 274\\
ICUR IV 11114 → 290\\
ICUR IV 11117a → 313\\
ICUR IV 11117b → 314\\
ICUR IV 11118 → 315\\
ICUR IV 11119 → 316\\
ICUR IV 11120 → 317\\
ICUR IV 11121 → 341\\
ICUR IV 11122 → 342\\
ICUR IV 11123 → 359\\
ICUR IV 11124a → 360\\
ICUR IV 11124d → 437\\
ICUR IV 11125 → 367\\
ICUR IV 11126a → 393\\
ICUR IV 11126b → 394\\
ICUR IV 11127 → 384\\
ICUR IV 11128 → 414\\
ICUR IV 11129 → 426\\
ICUR IV 11130 → 438\\
ICUR IV 11133 → 458\\
ICUR IV 11134 → 459\\
ICUR IV 11136 → 477\\
ICUR IV 11613t → 165\\
ICUR IV 11667u → 342\\
ICUR IV 11752 → 10\\
ICUR IV 11753 → 14\\
ICUR IV 11754 → 18\\
ICUR IV 11755 → 28\\
ICUR IV 11756 → 29\\
ICUR IV 11757 → 55\\
ICUR IV 11758 → 72\\
ICUR IV 11759 → 82\\
ICUR IV 11760 → 91\\
ICUR IV 11763 → 105\\
ICUR IV 11765 → 154\\
ICUR IV 11766 → 171\\
ICUR IV 11767 → 187\\
ICUR IV 11769 → 193\\
ICUR IV 11770 → 206\\
ICUR IV 11771 → 208\\
ICUR IV 11772 → 209\\
ICUR IV 11774 → 218\\
ICUR IV 11778 → 415\\
ICUR IV 11779 → 416\\
ICUR IV 12243 → 178\\
ICUR IV 12244 → 194\\
ICUR IV 12245 → 257\\
ICUR IV 12246 → 258\\
ICUR IV 12247 → 281\\
ICUR IV 12250 → 303\\
ICUR IV 12251 → 318\\
ICUR IV 12252a → 395\\
ICUR IV 12253 → 468\\
ICUR IV 12259a → 240\\
ICUR IV 12259e → 227\\
ICUR IV 12420 → 53\\
ICUR IV 12421 → 142\\
ICUR IV 12422 → 361\\
ICUR IV 12424 → 427\\
ICUR IV 12425 → 478\\
ICUR IV 12524 → 33\\
ICUR IV 12526 → 62\\
ICUR IV 12531 → 188\\
ICUR IV 12534 → 259\\
ICUR IV 12543 → 428\\
ICUR V 13105 → 56\\
ICUR V 13109 → 179\\
ICUR V 13355 → 306\\
ICUR VI 17248 → 230\\
ICUR VII 17437 → 50\\
ICUR VII 17438 → 49\\
ICUR VII 17475 → 216\\
ICUR VII 17494 → 295\\
ICUR VII 17495 → 296\\
ICUR VII 17510 → 386\\
ICUR VII 19963 → 260\\
ICUR VII 19974 → 330\\
ICUR VIII 20717 → 249\\
ICUR VIII 20718 → 297\\
ICUR VIII 20794 → 189\\
ICUR VIII 23061 → 144\\
ICUR VIII 23434 → 469\\
ICUR X 26312 → 20\\
ICUR X 26644 → 246\\
ICUR X 26645 → 286\\
ICUR X 27319 → 353\\
ICUR X 27324 → 397\end{multicols}


\section{Epigraphik-Datenbank Clauss/Slaby}

% généré par scripts/04_generated.py
\begin{multicols}{4}\footnotesize\noindent EDCS-01000473 → 306\\
EDCS-01000568 → 20\\
EDCS-02700238 → 396\\
EDCS-02700284 → 196\\
EDCS-03900021 → 34\\
EDCS-03900101 → 295\\
EDCS-03900102 → 386\\
EDCS-03900106 → 260\\
EDCS-05801896 → 439\\
EDCS-07000089 → 351\\
EDCS-07400001 → 329\\
EDCS-08300019 → 305\\
EDCS-08400008 → 417\\
EDCS-08900756 → 304\\
EDCS-09400838 → 195\\
EDCS-11700571 → 64\\
EDCS-11700572 → 78\\
EDCS-11700573 → 79\\
EDCS-11700574 → 128\\
EDCS-11700575 → 146\\
EDCS-11700578 → 405\\
EDCS-13000396 → 362\\
EDCS-13700250 → 241\\
EDCS-16300190 → 98\\
EDCS-16600245 → 282\\
EDCS-16600352 → 106\\
EDCS-16600356 → 228\\
EDCS-16600357 → 250\\
EDCS-18100517 → 36\\
EDCS-18200424 → 429\\
EDCS-18600145 → 62\\
EDCS-19000686 → 19\\
EDCS-19301140 → 296\\
EDCS-19400473 → 229\\
EDCS-19900369 → 213\\
EDCS-19900403 → 330\\
EDCS-19900430 → 179\\
EDCS-19900450 → 430\\
EDCS-19900458 → 262\\
EDCS-20000191 → 15\\
EDCS-20100025 → 321\\
EDCS-22300012 → 99\\
EDCS-23102819 → 397\\
EDCS-23201556 → 469\\
EDCS-23702281 → 368\\
EDCS-23702282 → 448\\
EDCS-24100416 → 189\\
EDCS-24100635 → 22\\
EDCS-25401950 → 348\\
EDCS-25402137 → 8\\
EDCS-25402147 → 447\\
EDCS-29900387 → 132\\
EDCS-30300391 → 353\\
EDCS-30600273 → 246\\
EDCS-30600274 → 286\\
EDCS-31900329 → 43\\
EDCS-31900337 → 382\\
EDCS-32900465 → 230\\
EDCS-33100921 → 56\\
EDCS-33500350 → 249\\
EDCS-33500351 → 297\\
EDCS-33900077 → 133\\
EDCS-33900078 → 234\\
EDCS-33900079 → 333\\
EDCS-33900148 → 38\\
EDCS-33900149 → 474\\
EDCS-33900175 → 42\\
EDCS-33900176 → 102\\
EDCS-33900177 → 347\\
EDCS-33900178 → 475\\
EDCS-33900327 → 101\\
EDCS-34100543 → 270\\
EDCS-34100544 → 277\\
EDCS-34100550 → 363\\
EDCS-34100563 → 129\\
EDCS-34100564 → 140\\
EDCS-34100565 → 148\\
EDCS-34100566 → 364\\
EDCS-34100568 → 420\\
EDCS-34100569 → 451\\
EDCS-34100570 → 476\\
EDCS-34100573 → 373\\
EDCS-34200648 → 52\\
EDCS-34200649 → 61\\
EDCS-34200650 → 87\\
EDCS-34200651 → 95\\
EDCS-34200652 → 118\\
EDCS-34200653 → 119\\
EDCS-34200654 → 120\\
EDCS-34200656 → 141\\
EDCS-34200657 → 151\\
EDCS-34200658 → 184\\
EDCS-34200659 → 217\\
EDCS-34200660 → 223\\
EDCS-34200663 → 278\\
EDCS-34200665 → 251\\
EDCS-34200666 → 272\\
EDCS-34200667 → 279\\
EDCS-34200668 → 287\\
EDCS-34200669 → 288\\
EDCS-34200670 → 299\\
EDCS-34200671 → 300\\
EDCS-34200672 → 301\\
EDCS-34200673 → 322\\
EDCS-34200674 → 323\\
EDCS-34200675 → 324\\
EDCS-34200676 → 336\\
EDCS-34200677 → 337\\
EDCS-34200678 → 349\\
EDCS-34200679 → 376\\
EDCS-34200681 → 408\\
EDCS-34200682 → 409\\
EDCS-34200683 → 410\\
EDCS-34200684 → 411\\
EDCS-34200685 → 443\\
EDCS-34200686 → 444\\
EDCS-34200687 → 453\\
EDCS-34200688 → 454\\
EDCS-34200689 → 455\\
EDCS-34200690 → 466\\
EDCS-34201004 → 160\\
EDCS-34201006 → 271\\
EDCS-34201007 → 334\\
EDCS-34300812 → 365\\
EDCS-34300830 → 434\\
EDCS-34501139 → 112\\
EDCS-34501143 → 4\\
EDCS-34501146 → 244\\
EDCS-34501147 → 275\\
EDCS-34700228 → 110\\
EDCS-34700234 → 168\\
EDCS-34801295 → 21\\
EDCS-34801298 → 27\\
EDCS-34801299 → 46\\
EDCS-34801300 → 66\\
EDCS-34801302 → 121\\
EDCS-34801303 → 162\\
EDCS-34801304 → 237\\
EDCS-34801306 → 311\\
EDCS-34801307 → 312\\
EDCS-34801308 → 350\\
EDCS-34801309 → 377\\
EDCS-34801310 → 392\\
EDCS-34801311 → 435\\
EDCS-34900709 → 88\\
EDCS-34900710 → 96\\
EDCS-34900711 → 252\\
EDCS-35100216 → 39\\
EDCS-35100217 → 57\\
EDCS-35100218 → 59\\
EDCS-35100219 → 67\\
EDCS-35100466 → 470\\
EDCS-35100467 → 352\\
EDCS-35100468 → 433\\
EDCS-35200265 → 49\\
EDCS-35200283 → 214\\
EDCS-35200401 → 73\\
EDCS-35200402 → 125\\
EDCS-35200404 → 440\\
EDCS-35300743 → 144\\
EDCS-35700175 → 122\\
EDCS-35700176 → 131\\
EDCS-35700177 → 134\\
EDCS-35700179 → 238\\
EDCS-35700180 → 239\\
EDCS-35700181 → 253\\
EDCS-35700183 → 289\\
EDCS-35700184 → 294\\
EDCS-35700185 → 325\\
EDCS-35700186 → 338\\
EDCS-35700187 → 339\\
EDCS-35700188 → 355\\
EDCS-35700189 → 356\\
EDCS-35700190 → 378\\
EDCS-35700191 → 379\\
EDCS-35700192 → 380\\
EDCS-35700193 → 381\\
EDCS-35700194 → 421\\
EDCS-35700195 → 436\\
EDCS-35700196 → 456\\
EDCS-35700197 → 445\\
EDCS-35700198 → 457\\
EDCS-35700212 → 422\\
EDCS-35700729 → 395\\
EDCS-35700730 → 240\\
EDCS-35700732 → 227\\
EDCS-35700797 → 174\\
EDCS-35700798 → 175\\
EDCS-35700799 → 192\\
EDCS-35700937 → 18\\
EDCS-35700942 → 105\\
EDCS-35700944 → 154\\
EDCS-35700945 → 171\\
EDCS-35700946 → 206\\
EDCS-35802005 → 2\\
EDCS-35802006 → 30\\
EDCS-35802007 → 31\\
EDCS-35802008 → 68\\
EDCS-35802010 → 204\\
EDCS-35802012 → 254\\
EDCS-35802013 → 273\\
EDCS-35802015 → 340\\
EDCS-35802018 → 358\\
EDCS-35802019 → 423\\
EDCS-35802020 → 467\\
EDCS-35900808 → 326\\
EDCS-35900809 → 123\\
EDCS-35900810 → 170\\
EDCS-35900834 → 449\\
EDCS-35900846 → 424\\
EDCS-35900847 → 366\\
EDCS-35900965 → 9\\
EDCS-35900967 → 48\\
EDCS-35900968 → 169\\
EDCS-35900969 → 185\\
EDCS-35900970 → 186\\
EDCS-35900972 → 255\\
EDCS-35900973 → 446\\
EDCS-36100891 → 383\\
EDCS-36101085 → 6\\
EDCS-36101086 → 69\\
EDCS-36101087 → 70\\
EDCS-36101088 → 81\\
EDCS-36101089 → 89\\
EDCS-36101090 → 97\\
EDCS-36101092 → 124\\
EDCS-36101094 → 163\\
EDCS-36101095 → 191\\
EDCS-36101096 → 207\\
EDCS-36101097 → 224\\
EDCS-36101098 → 225\\
EDCS-36101099 → 280\\
EDCS-36101100 → 302\\
EDCS-36101101 → 327\\
EDCS-36101103 → 328\\
EDCS-36101106 → 412\\
EDCS-36101107 → 413\\
EDCS-36101108 → 425\\
EDCS-36101279 → 12\\
EDCS-36101280 → 13\\
EDCS-36101281 → 25\\
EDCS-36101283 → 45\\
EDCS-36101285 → 65\\
EDCS-36101286 → 116\\
EDCS-36101288 → 117\\
EDCS-36101289 → 161\\
EDCS-36101297 → 202\\
EDCS-36101305 → 310\\
EDCS-36101309 → 374\\
EDCS-36101310 → 406\\
EDCS-36101312 → 407\\
EDCS-36101314 → 452\\
EDCS-36901720 → 103\\
EDCS-37200371 → 3\\
EDCS-37200372 → 32\\
EDCS-37200373 → 44\\
EDCS-37200375 → 58\\
EDCS-37200376 → 71\\
EDCS-37200378 → 90\\
EDCS-37200379 → 104\\
EDCS-37200383 → 135\\
EDCS-37200384 → 152\\
EDCS-37200385 → 153\\
EDCS-37200386 → 164\\
EDCS-37200387 → 176\\
EDCS-37200388 → 177\\
EDCS-37200389 → 205\\
EDCS-37200391 → 226\\
EDCS-37200393 → 256\\
EDCS-37200394 → 274\\
EDCS-37200396 → 290\\
EDCS-37200399 → 313\\
EDCS-37200400 → 314\\
EDCS-37200401 → 315\\
EDCS-37200402 → 316\\
EDCS-37200403 → 317\\
EDCS-37200404 → 341\\
EDCS-37200405 → 342\\
EDCS-37200406 → 359\\
EDCS-37200407 → 360\\
EDCS-37200410 → 437\\
EDCS-37200411 → 367\\
EDCS-37200412 → 393\\
EDCS-37200413 → 394\\
EDCS-37200414 → 384\\
EDCS-37200415 → 414\\
EDCS-37200416 → 426\\
EDCS-37200417 → 438\\
EDCS-37200420 → 458\\
EDCS-37200421 → 459\\
EDCS-37200423 → 477\\
EDCS-37300523 → 165\\
EDCS-37400847 → 10\\
EDCS-37400848 → 14\\
EDCS-37400849 → 28\\
EDCS-37400850 → 29\\
EDCS-37400851 → 55\\
EDCS-37400852 → 72\\
EDCS-37400853 → 82\\
EDCS-37400854 → 91\\
EDCS-37400856 → 187\\
EDCS-37400858 → 193\\
EDCS-37400859 → 208\\
EDCS-37400860 → 209\\
EDCS-37400862 → 218\\
EDCS-37400866 → 415\\
EDCS-37400867 → 416\\
EDCS-37600678 → 178\\
EDCS-37600679 → 194\\
EDCS-37600680 → 257\\
EDCS-37600681 → 258\\
EDCS-37600682 → 281\\
EDCS-37600684 → 303\\
EDCS-37600685 → 318\\
EDCS-37600687 → 468\\
EDCS-37600944 → 53\\
EDCS-37600945 → 142\\
EDCS-37600946 → 361\\
EDCS-37600947 → 427\\
EDCS-37600948 → 478\\
EDCS-37700147 → 212\\
EDCS-37700196 → 183\\
EDCS-37801588 → 33\\
EDCS-37801594 → 188\\
EDCS-37801595 → 259\\
EDCS-37801603 → 428\\
EDCS-38100789 → 130\\
EDCS-38100790 → 150\\
EDCS-38100791 → 309\\
EDCS-38100792 → 335\\
EDCS-38100823 → 236\\
EDCS-38100824 → 298\\
EDCS-38302409 → 26\\
EDCS-38302410 → 149\\
EDCS-38302491 → 86\\
EDCS-38302492 → 203\\
EDCS-38302493 → 235\\
EDCS-38302567 → 375\\
EDCS-38600081 → 263\\
EDCS-38600097 → 23\\
EDCS-38600098 → 63\\
EDCS-38600099 → 155\\
EDCS-38600100 → 197\\
EDCS-38600196 → 264\\
EDCS-38600270 → 7\\
EDCS-38600271 → 108\\
EDCS-38600308 → 431\\
EDCS-38600315 → 24\\
EDCS-38600325 → 265\\
EDCS-38600338 → 450\\
EDCS-38600420 → 137\\
EDCS-38700412 → 109\\
EDCS-38700455 → 400\\
EDCS-38700473 → 74\\
EDCS-38700493 → 343\\
EDCS-38700644 → 11\\
EDCS-38700710 → 243\\
EDCS-38700730 → 92\\
EDCS-38700748 → 219\\
EDCS-38700749 → 267\\
EDCS-38700751 → 370\\
EDCS-38700752 → 461\\
EDCS-38700879 → 398\\
EDCS-38700887 → 126\\
EDCS-38700889 → 266\\
EDCS-38700907 → 242\\
EDCS-38700928 → 432\\
EDCS-38700931 → 441\\
EDCS-38701118 → 399\\
EDCS-38701119 → 460\\
EDCS-38701126 → 127\\
EDCS-38801380 → 50\\
EDCS-38801427 → 216\\
EDCS-39000328 → 111\\
EDCS-39000329 → 138\\
EDCS-39000330 → 220\\
EDCS-39000331 → 284\\
EDCS-39000332 → 344\\
EDCS-39000407 → 1\\
EDCS-39000408 → 17\\
EDCS-39000409 → 35\\
EDCS-39000410 → 54\\
EDCS-39000411 → 75\\
EDCS-39000413 → 100\\
EDCS-39000414 → 113\\
EDCS-39000415 → 143\\
EDCS-39000416 → 156\\
EDCS-39000417 → 157\\
EDCS-39000418 → 198\\
EDCS-39000421 → 215\\
EDCS-39000422 → 231\\
EDCS-39000424 → 285\\
EDCS-39000425 → 292\\
EDCS-39000426 → 319\\
EDCS-39000429 → 388\\
EDCS-39000430 → 401\\
EDCS-39000431 → 418\\
EDCS-39000433 → 471\\
EDCS-39000434 → 472\\
EDCS-39101961 → 5\\
EDCS-39101962 → 40\\
EDCS-39101964 → 166\\
EDCS-39101966 → 462\\
EDCS-39101967 → 283\\
EDCS-39101968 → 291\\
EDCS-39200296 → 180\\
EDCS-39200298 → 51\\
EDCS-39200381 → 172\\
EDCS-39200382 → 173\\
EDCS-39200384 → 354\\
EDCS-39200385 → 419\\
EDCS-39200386 → 463\\
EDCS-39300305 → 37\\
EDCS-39300306 → 76\\
EDCS-39300307 → 77\\
EDCS-39300308 → 84\\
EDCS-39300309 → 93\\
EDCS-39300310 → 114\\
EDCS-39300312 → 145\\
EDCS-39300313 → 158\\
EDCS-39300314 → 167\\
EDCS-39300315 → 181\\
EDCS-39300316 → 199\\
EDCS-39300317 → 221\\
EDCS-39300318 → 245\\
EDCS-39300319 → 268\\
EDCS-39300321 → 320\\
EDCS-39300322 → 331\\
EDCS-39300323 → 345\\
EDCS-39300324 → 371\\
EDCS-39300325 → 389\\
EDCS-39300326 → 390\\
EDCS-39300327 → 402\\
EDCS-39300328 → 403\\
EDCS-39300329 → 442\\
EDCS-39900259 → 94\\
EDCS-39900318 → 41\\
EDCS-39900319 → 85\\
EDCS-39900320 → 308\\
EDCS-39900321 → 404\\
EDCS-39900335 → 107\\
EDCS-39900394 → 247\\
EDCS-39900423 → 60\\
EDCS-39900424 → 346\\
EDCS-39900484 → 248\\
EDCS-39900588 → 16\\
EDCS-39900589 → 232\\
EDCS-39900598 → 210\\
EDCS-40300313 → 47\\
EDCS-40300314 → 80\\
EDCS-40300321 → 115\\
EDCS-40300325 → 136\\
EDCS-40300326 → 139\\
EDCS-40300327 → 147\\
EDCS-40300329 → 159\\
EDCS-40300330 → 182\\
EDCS-40300332 → 190\\
EDCS-40300334 → 200\\
EDCS-40300335 → 201\\
EDCS-40300336 → 211\\
EDCS-40300339 → 222\\
EDCS-40300340 → 233\\
EDCS-40300341 → 269\\
EDCS-40300342 → 261\\
EDCS-40300343 → 276\\
EDCS-40300344 → 293\\
EDCS-40300347 → 332\\
EDCS-40300349 → 372\\
EDCS-40300353 → 464\\
EDCS-40300354 → 473\\
EDCS-40301026 → 465\\
EDCS-40301037 → 391\\
EDCS-43700747 → 357\\
EDCS-44100241 → 369\\
EDCS-59400075 → 83\\
EDCS-71500011 → 387\\
EDCS-74700288 → 307\\
EDCS-75200321 → 385\end{multicols}


\section{Epigraphic Database Bari}

% généré par scripts/04_generated.py
\begin{multicols}{4}\footnotesize\noindent EDB699 → 10\\
EDB700 → 14\\
EDB746 → 360\\
EDB876 → 187\\
EDB878 → 193\\
EDB880 → 206\\
EDB965 → 18\\
EDB966 → 28\\
EDB1004 → 105\\
EDB1342 → 208\\
EDB2091 → 82\\
EDB2092 → 91\\
EDB2099 → 154\\
EDB2100 → 171\\
EDB3284 → 228\\
EDB3285 → 278\\
EDB4428 → 56\\
EDB4650 → 240\\
EDB4654 → 227\\
EDB5083 → 194\\
EDB5084 → 257\\
EDB5167 → 415\\
EDB5168 → 416\\
EDB5268 → 209\\
EDB5272 → 218\\
EDB5345 → 258\\
EDB5346 → 281\\
EDB5349 → 303\\
EDB5945 → 179\\
EDB6536 → 318\\
EDB6728 → 29\\
EDB6729 → 55\\
EDB6730 → 72\\
EDB7067 → 395\\
EDB7068 → 468\\
EDB8228 → 42\\
EDB8229 → 102\\
EDB8750 → 250\\
EDB8751 → 251\\
EDB8752 → 272\\
EDB8753 → 279\\
EDB8839 → 178\\
EDB9633 → 141\\
EDB10125 → 324\\
EDB10126 → 329\\
EDB10126 → 336\\
EDB10655 → 20\\
EDB11272 → 348\\
EDB11273 → 475\\
EDB11277 → 287\\
EDB11278 → 282\\
EDB11278 → 288\\
EDB11279 → 299\\
EDB11280 → 300\\
EDB11281 → 301\\
EDB11282 → 322\\
EDB11297 → 168\\
EDB11967 → 347\\
EDB12186 → 21\\
EDB12500 → 447\\
EDB12805 → 230\\
EDB13166 → 119\\
EDB13167 → 106\\
EDB13167 → 120\\
EDB13470 → 160\\
EDB13472 → 271\\
EDB13473 → 334\\
EDB13893 → 453\\
EDB13894 → 454\\
EDB13895 → 455\\
EDB13896 → 466\\
EDB14639 → 323\\
EDB14891 → 337\\
EDB14892 → 349\\
EDB14893 → 376\\
EDB15528 → 189\\
EDB15588 → 8\\
EDB16486 → 408\\
EDB16487 → 409\\
EDB16488 → 410\\
EDB16489 → 411\\
EDB16490 → 443\\
EDB16491 → 444\\
EDB16941 → 234\\
EDB16942 → 262\\
EDB16943 → 333\\
EDB17188 → 302\\
EDB17189 → 327\\
EDB17191 → 328\\
EDB17432 → 151\\
EDB17433 → 184\\
EDB17434 → 217\\
EDB17435 → 223\\
EDB17462 → 52\\
EDB17463 → 61\\
EDB17464 → 87\\
EDB17465 → 95\\
EDB17466 → 118\\
EDB17657 → 38\\
EDB17710 → 216\\
EDB17954 → 386\\
EDB17958 → 192\\
EDB17959 → 207\\
EDB17960 → 224\\
EDB17961 → 225\\
EDB17962 → 280\\
EDB18146 → 49\\
EDB18365 → 97\\
EDB18367 → 124\\
EDB18369 → 163\\
EDB18370 → 174\\
EDB18371 → 175\\
EDB18372 → 191\\
EDB18727 → 70\\
EDB18778 → 9\\
EDB18940 → 373\\
EDB19129 → 424\\
EDB19138 → 270\\
EDB19139 → 277\\
EDB19166 → 213\\
EDB19222 → 101\\
EDB19224 → 133\\
EDB19258 → 81\\
EDB19321 → 413\\
EDB19322 → 425\\
EDB19494 → 366\\
EDB19574 → 383\\
EDB19681 → 110\\
EDB19792 → 420\\
EDB19793 → 451\\
EDB19794 → 476\\
EDB19824 → 89\\
EDB19876 → 46\\
EDB19877 → 66\\
EDB19879 → 121\\
EDB20076 → 98\\
EDB20295 → 330\\
EDB20348 → 50\\
EDB20435 → 449\\
EDB20443 → 412\\
EDB20583 → 255\\
EDB20584 → 382\\
EDB20586 → 446\\
EDB20721 → 27\\
EDB20813 → 236\\
EDB20818 → 365\\
EDB20872 → 6\\
EDB20873 → 69\\
EDB20916 → 129\\
EDB20917 → 140\\
EDB20918 → 148\\
EDB20919 → 364\\
EDB20990 → 48\\
EDB20991 → 99\\
EDB20992 → 169\\
EDB20993 → 185\\
EDB20994 → 186\\
EDB21154 → 130\\
EDB21298 → 298\\
EDB21312 → 183\\
EDB21373 → 295\\
EDB21374 → 296\\
EDB21458 → 150\\
EDB21459 → 309\\
EDB21460 → 335\\
EDB21504 → 260\\
EDB21579 → 469\\
EDB22029 → 375\\
EDB22056 → 235\\
EDB22061 → 162\\
EDB22062 → 237\\
EDB22064 → 311\\
EDB22065 → 312\\
EDB22066 → 350\\
EDB22067 → 377\\
EDB22177 → 203\\
EDB22178 → 86\\
EDB22305 → 392\\
EDB22306 → 435\\
EDB24129 → 88\\
EDB24130 → 96\\
EDB24131 → 252\\
EDB24663 → 39\\
EDB24664 → 57\\
EDB24665 → 59\\
EDB24666 → 67\\
EDB24669 → 122\\
EDB24670 → 131\\
EDB24671 → 134\\
EDB24673 → 238\\
EDB24674 → 239\\
EDB24675 → 253\\
EDB24676 → 289\\
EDB24677 → 294\\
EDB24678 → 325\\
EDB24679 → 338\\
EDB24680 → 339\\
EDB24681 → 355\\
EDB24682 → 356\\
EDB24684 → 379\\
EDB24685 → 380\\
EDB24686 → 381\\
EDB24687 → 421\\
EDB24688 → 436\\
EDB24689 → 456\\
EDB24690 → 445\\
EDB24691 → 457\\
EDB24704 → 422\\
EDB24920 → 326\\
EDB24921 → 123\\
EDB24922 → 170\\
EDB24939 → 358\\
EDB24940 → 423\\
EDB24941 → 467\\
EDB25348 → 378\\
EDB25692 → 2\\
EDB25693 → 30\\
EDB25694 → 31\\
EDB25695 → 43\\
EDB25696 → 68\\
EDB25698 → 204\\
EDB25700 → 254\\
EDB25701 → 273\\
EDB25702 → 340\\
EDB25703 → 357\\
EDB27045 → 306\\
EDB27137 → 36\\
EDB27212 → 15\\
EDB27361 → 304\\
EDB27435 → 263\\
EDB27454 → 23\\
EDB27455 → 63\\
EDB27456 → 155\\
EDB27457 → 197\\
EDB27580 → 264\\
EDB27782 → 7\\
EDB27783 → 108\\
EDB27849 → 431\\
EDB27940 → 24\\
EDB27988 → 265\\
EDB28045 → 450\\
EDB28273 → 137\\
EDB28397 → 398\\
EDB28428 → 126\\
EDB28430 → 266\\
EDB28484 → 242\\
EDB28503 → 432\\
EDB28506 → 441\\
EDB28939 → 470\\
EDB28944 → 399\\
EDB28945 → 460\\
EDB28952 → 127\\
EDB28954 → 352\\
EDB29417 → 22\\
EDB29441 → 109\\
EDB29526 → 400\\
EDB29591 → 125\\
EDB29599 → 74\\
EDB29619 → 343\\
EDB29620 → 430\\
EDB30041 → 11\\
EDB30112 → 243\\
EDB30135 → 92\\
EDB30156 → 219\\
EDB30157 → 267\\
EDB30159 → 370\\
EDB30165 → 461\\
EDB30459 → 5\\
EDB30464 → 40\\
EDB30476 → 166\\
EDB30485 → 283\\
EDB30488 → 291\\
EDB30676 → 111\\
EDB30683 → 138\\
EDB30687 → 220\\
EDB30696 → 284\\
EDB30697 → 344\\
EDB30801 → 112\\
EDB31183 → 1\\
EDB31184 → 4\\
EDB31185 → 17\\
EDB31186 → 35\\
EDB31187 → 54\\
EDB31188 → 75\\
EDB31189 → 73\\
EDB31193 → 100\\
EDB31195 → 113\\
EDB31196 → 143\\
EDB31197 → 156\\
EDB31198 → 157\\
EDB31199 → 198\\
EDB31203 → 215\\
EDB31204 → 231\\
EDB31205 → 244\\
EDB31206 → 275\\
EDB31208 → 292\\
EDB31209 → 319\\
EDB31210 → 285\\
EDB31230 → 307\\
EDB31239 → 388\\
EDB31240 → 401\\
EDB31242 → 418\\
EDB31243 → 433\\
EDB31244 → 440\\
EDB31246 → 471\\
EDB31247 → 472\\
EDB31761 → 180\\
EDB31763 → 51\\
EDB31779 → 363\\
EDB31855 → 172\\
EDB31856 → 173\\
EDB31858 → 354\\
EDB31859 → 419\\
EDB31861 → 463\\
EDB32043 → 37\\
EDB32044 → 76\\
EDB32045 → 77\\
EDB32048 → 84\\
EDB32049 → 93\\
EDB32051 → 114\\
EDB32053 → 145\\
EDB32056 → 158\\
EDB32057 → 167\\
EDB32058 → 181\\
EDB32059 → 199\\
EDB32075 → 221\\
EDB32081 → 229\\
EDB32082 → 245\\
EDB32083 → 268\\
EDB32090 → 320\\
EDB32097 → 331\\
EDB32098 → 345\\
EDB32099 → 371\\
EDB32100 → 389\\
EDB32101 → 390\\
EDB32105 → 402\\
EDB32106 → 403\\
EDB32110 → 442\\
EDB32543 → 34\\
EDB32628 → 94\\
EDB32698 → 41\\
EDB32699 → 85\\
EDB32700 → 308\\
EDB32701 → 404\\
EDB32714 → 107\\
EDB32728 → 64\\
EDB32729 → 78\\
EDB32730 → 79\\
EDB32731 → 128\\
EDB32732 → 146\\
EDB32736 → 405\\
EDB32852 → 247\\
EDB32878 → 60\\
EDB32879 → 346\\
EDB32975 → 248\\
EDB33080 → 16\\
EDB33081 → 232\\
EDB33094 → 210\\
EDB33117 → 47\\
EDB33118 → 80\\
EDB33127 → 115\\
EDB33131 → 136\\
EDB33132 → 139\\
EDB33133 → 147\\
EDB33135 → 159\\
EDB33136 → 182\\
EDB33138 → 190\\
EDB33139 → 200\\
EDB33140 → 201\\
EDB33143 → 211\\
EDB33145 → 214\\
EDB33147 → 222\\
EDB33148 → 233\\
EDB33150 → 269\\
EDB33151 → 261\\
EDB33152 → 276\\
EDB33154 → 293\\
EDB33157 → 332\\
EDB33161 → 372\\
EDB33162 → 439\\
EDB33166 → 464\\
EDB33167 → 473\\
EDB33615 → 195\\
EDB34225 → 465\\
EDB34240 → 391\\
EDB34442 → 305\\
EDB34468 → 434\\
EDB35329 → 474\\
EDB35372 → 462\\
EDB35388 → 353\\
EDB35584 → 144\\
EDB35608 → 246\\
EDB35609 → 286\\
EDB35683 → 249\\
EDB35685 → 297\\
EDB35690 → 397\\
EDB38747 → 417\\
EDB39161 → 19\\
EDB39162 → 103\\
EDB39548 → 3\\
EDB39549 → 32\\
EDB39550 → 44\\
EDB39552 → 58\\
EDB39553 → 71\\
EDB39555 → 90\\
EDB39556 → 104\\
EDB39560 → 135\\
EDB39561 → 152\\
EDB39562 → 153\\
EDB39563 → 164\\
EDB39564 → 176\\
EDB39565 → 177\\
EDB39569 → 205\\
EDB39571 → 226\\
EDB39573 → 256\\
EDB39574 → 274\\
EDB39576 → 290\\
EDB39579 → 313\\
EDB39579 → 314\\
EDB39580 → 315\\
EDB39581 → 316\\
EDB39582 → 317\\
EDB39583 → 341\\
EDB39584 → 342\\
EDB39585 → 359\\
EDB39588 → 437\\
EDB39589 → 367\\
EDB39590 → 393\\
EDB39590 → 394\\
EDB39591 → 384\\
EDB39592 → 414\\
EDB39593 → 426\\
EDB39594 → 438\\
EDB39597 → 458\\
EDB39598 → 459\\
EDB39600 → 477\\
EDB40555 → 165\\
EDB41240 → 351\\
EDB41423 → 62\\
EDB41696 → 142\\
EDB41699 → 361\\
EDB41702 → 427\\
EDB41720 → 53\\
EDB41721 → 478\\
EDB41825 → 428\\
EDB41941 → 368\\
EDB41942 → 448\\
EDB42064 → 369\\
EDB42430 → 259\\
EDB42433 → 188\\
EDB42440 → 33\\
EDB42643 → 321\\
EDB42676 → 396\\
EDB42969 → 241\\
EDB43052 → 310\\
EDB43058 → 362\end{multicols}


\section{Autres recueils et publications}

Références au \emph{Corpus Inscriptionum Latinarum}, aux \emph{Inscriptiones Latinae Christianae Veteres}, aux \emph{Carmina Latina Epigraphica}, à \emph{L’Année épigraphique} et aux autres publications, dans la forme abrégée des bases.

% généré par scripts/04_generated.py
\begin{multicols}{3}\footnotesize\noindent AE 1892, 31 → 126\\
AE 1904, 33 → 72\\
AE 1904, 122 → 72\\
AE 1905, 71 → 106\\
AE 1905, 72 → 119\\
AE 1905, 75 → 228\\
AE 1905, 76 → 250\\
AE 1905, 77 → 251\\
AE 1905, 78 → 272\\
AE 1905, 79 → 287\\
AE 1905, 80 → 300\\
AE 1905, 81 → 324\\
AE 1905, 82 → 408\\
AE 1905, 83 → 444\\
AE 1905, 84 → 454\\
AE 1905, 93 → 466\\
AE 1906, 82 → 376\\
AE 1906, 83 → 61\\
AE 1906, 85 → 323\\
AE 1906, 86 → 282\\
AE 1907, 133 → 15\\
AE 1909, 27 → 397\\
AE 1910, 205 → 447\\
AE 1912, 261 → 134\\
AE 1914, 68 → 98\\
AE 1933, 158 → 397\\
AE 1939, 152 → 230\\
AE 1949, 163 → 230\\
AE 1959, 64 → 321\\
AE 1974, 85 → 195\\
AE 1979, 43 → 304\\
AE 1984, 51 → 417\\
AE 1985, 60 → 305\\
AE 1987, 61 → 329\\
AE 1987, 61 → 336\\
AE 1988, 34 → 20\\
AE 1988, 111 → 351\\
AE 1994, 163 → 396\\
AE 1994, 263 → 196\\
AE 1997, 166 → 4\\
AE 1997, 166 → 229\\
AE 1997, 166 → 244\\
AE 1997, 166 → 275\\
AE 1997, 166 → 296\\
AE 1999, 307 → 241\\
AE 2000, 187 → 321\\
AE 2001, 219 → 396\\
AE 2001, 526 → 368\\
AE 2001, 527 → 448\\
AE 2003, 202 → 321\\
AE 2003, 216 → 132\\
AE 2006, 93 → 36\\
AE 2006, 240 → 369\\
AE 2009, 140 → 321\\
AE 2009, 145 → 179\\
AE 2009, 145 → 306\\
AE 2010, 207 → 83\\
AE 2012, 197 → 26\\
AE 2016, 170 → 385\\
AE 2017, 91 → 428\\
AIIRoma-05, p 111 → 444\\
Anagni 00041 → 34\\
Anagni 00150 → 295\\
Anagni 00151 → 386\\
Anagni 00155 → 260\\
AnalEpi p 69 → 230\\
AnalEpi p 101 → 134\\
AnalEpi p 377 → 319\\
Arctos-2012-225 → 26\\
BE 00663 → 475\\
Bagno 00051 → 94\\
Bencivenni-2018, 38,1 → 60\\
CEACelio 00311 → 362\\
CEACelio 00466 → 396\\
CIL 5 429 → 210\\
CIL 5 672,67 → 261\\
CIL 9 630,3 → 449\\
CIL 10 949,05 → 16\\
CIL 10 951,6 → 49\\
CIL 10 1088,458 → 99\\
CIL 13 381 → 136\\
CIL VI 1712 (p 3173, 4741) → 36\\
CIL VI 7429 → 62\\
CIL VI 8405 → 429\\
CIL VI 9237 → 19\\
CIL VI 9597 (p 3470) → 296\\
CIL VI 9787 → 229\\
CIL VI 31972 (p 4798) → 213\\
CIL VI 32011 (p 4803) → 59\\
CIL VI 32018 (p 4804) → 330\\
CIL VI 32045a (p 4806) → 179\\
CIL VI 32069 (p 4809) → 430\\
CIL VI 32078 (p 4810) → 262\\
CIL VI 32955 → 469\\
CIL VI 33865 → 22\\
CIL VI 33902 → 189\\
CIL VI 37122 (p 4820) → 15\\
CIL VI 39147a → 397\\
CIL VI 39440b (p 4068) → 330\\
CIL VI 41342 → 306\\
CIL VI 41424a → 20\\
CIL XIV 290,12 → 107\\
CIL XIV 3896 → 439\\
CLE 659 → 22\\
CLE 660 → 43\\
CLE 662 → 73\\
CLE 664 → 125\\
CLE 670 → 244\\
CLE 676 → 382\\
CLE 679 → 433\\
CLE 680 → 440\\
CLE 682 → 470\\
CLE 1346 → 352\\
CLE 1619 → 214\\
CLE 2005 → 397\\
CLENuovo p 108 → 244\\
Carletti-2008, 56 → 70\\
Carletti-2008, 58 → 202\\
Carletti-2008, 95 → 324\\
Carletti-2008, 177 → 387\\
Collatia p 639 → 230\\
Damaso p 110 → 244\\
Epigraphica-1985-109 → 411\\
Epigraphica-1986-186 → 49\\
Epigraphica-2017-417 → 428\\
Ferrua-1979, 1444 → 307\\
Gummerus-01, 98 → 296\\
HEp 2006, 435 → 369\\
ICVaticano p 284 → 244\\
ICVaticano p 299 → 4\\
ICVaticano p 310 → 296\\
ICVaticano p 334 → 229\\
ICVaticano p 350 → 275\\
ICaRoma 00026 → 324\\
ICaRoma 00070 → 20\\
ICaRoma 00085 → 70\\
ICaRoma 00089 → 4\\
ICaRoma 00124 → 310\\
ICaRoma 00130 → 202\\
IGLFRPal 00201 → 112\\
ILCV 94a → 93\\
ILCV 95 → 179\\
ILCV 120 → 429\\
ILCV 139 → 213\\
ILCV 139 → 430\\
ILCV 162 (add) → 15\\
ILCV 172 → 262\\
ILCV 197 (add) → 330\\
ILCV 288 (add, em) → 59\\
ILCV 301 (add) → 388\\
ILCV 317 (add) → 440\\
ILCV 569 → 333\\
ILCV 589 → 19\\
ILCV 610 → 296\\
ILCV 657 (em) → 189\\
ILCV 669 → 229\\
ILCV 732 (add) → 469\\
ILCV 748 (add) → 22\\
ILCV 775 (em) → 183\\
ILCV 994 (em) → 378\\
ILCV 1002 → 203\\
ILCV 1002 → 400\\
ILCV 1119 → 289\\
ILCV 1126a (em) → 105\\
ILCV 1126b → 18\\
ILCV 1128 → 418\\
ILCV 1128 (add, em) → 134\\
ILCV 1130 → 310\\
ILCV 1268 (em) → 67\\
ILCV 1269 → 421\\
ILCV 1269 (em) → 202\\
ILCV 1270 (add) → 269\\
ILCV 1296a (add) → 111\\
ILCV 1297 (em) → 142\\
ILCV 1327 → 81\\
ILCV 1327a (em) → 94\\
ILCV 1328 (em) → 191\\
ILCV 1464 → 214\\
ILCV 1468 → 458\\
ILCV 1478a → 145\\
ILCV 1479 → 210\\
ILCV 1479a → 233\\
ILCV 1480 → 276\\
ILCV 1481 → 308\\
ILCV 1483 (add) → 391\\
ILCV 1508 → 405\\
ILCV 1548 → 201\\
ILCV 1563 → 70\\
ILCV 1599 → 352\\
ILCV 1603 (add, em) → 133\\
ILCV 1605 → 474\\
ILCV 1614 (em) → 152\\
ILCV 1735 → 45\\
ILCV 1850 → 36\\
ILCV 1857 (em) → 321\\
ILCV 2007 (add) → 172\\
ILCV 2090 → 57\\
ILCV 2132 → 339\\
ILCV 2146 (em) → 381\\
ILCV 2148 → 232\\
ILCV 2281 → 446\\
ILCV 2398 (addd) → 475\\
ILCV 2451 (em) → 43\\
ILCV 2569 → 61\\
ILCV 2570 → 419\\
ILCV 2571 (em) → 9\\
ILCV 2576 (add) → 463\\
ILCV 2602 (em) → 16\\
ILCV 2603 → 75\\
ILCV 2603a → 140\\
ILCV 2604 → 166\\
ILCV 2606 (em) → 376\\
ILCV 2609 (em) → 117\\
ILCV 2627 → 17\\
ILCV 2628 → 121\\
ILCV 2628 (em) → 35\\
ILCV 2633 (add) → 353\\
ILCV 2690 → 53\\
ILCV 2690 (em) → 157\\
ILCV 2732 → 125\\
ILCV 2772 → 211\\
ILCV 2782 (add) → 268\\
ILCV 2795b → 33\\
ILCV 2795c (add, em) → 300\\
ILCV 2807a → 58\\
ILCV 2807a → 116\\
ILCV 2816 (em) → 4\\
ILCV 2819 (em) → 9\\
ILCV 2819a (em) → 260\\
ILCV 2824 (add) → 118\\
ILCV 2846 → 6\\
ILCV 2851 → 223\\
ILCV 2868 → 250\\
ILCV 2874a → 144\\
ILCV 2878 → 110\\
ILCV 2879 (em) → 271\\
ILCV 2883 (add, em) → 56\\
ILCV 2925a → 299\\
ILCV 2926 → 95\\
ILCV 2927a (add, em) → 466\\
ILCV 2940 → 23\\
ILCV 2940a → 34\\
ILCV 2941 → 50\\
ILCV 2942 → 100\\
ILCV 2943 (em) → 109\\
ILCV 2945a → 208\\
ILCV 2946 → 386\\
ILCV 2946 → 425\\
ILCV 2967 → 66\\
ILCV 2969 (em) → 205\\
ILCV 2969a → 279\\
ILCV 2971 (em) → 336\\
ILCV 2971a → 367\\
ILCV 2974 (em) → 451\\
ILCV 2976a → 119\\
ILCV 2977 → 323\\
ILCV 2978a → 41\\
ILCV 2978b → 68\\
ILCV 2978c → 288\\
ILCV 3003a (em) → 383\\
ILCV 3057a (add) → 255\\
ILCV 3099 → 90\\
ILCV 3099 → 120\\
ILCV 3099 (em) → 42\\
ILCV 3099a → 371\\
ILCV 3099b → 404\\
ILCV 3137 (em) → 80\\
ILCV 3142a → 460\\
ILCV 3252a (em) → 38\\
ILCV 3252b → 99\\
ILCV 3253 → 124\\
ILCV 3253a (em) → 151\\
ILCV 3287a (em) → 311\\
ILCV 3305 → 182\\
ILCV 3306 (em) → 295\\
ILCV 3359 → 73\\
ILCV 3382 (em) → 27\\
ILCV 3390a → 249\\
ILCV 3424 → 382\\
ILCV 3432 → 433\\
ILCV 3444 → 400\\
ILCV 3446 (add) → 244\\
ILCV 3501 → 320\\
ILCV 3502 → 347\\
ILCV 3503 → 454\\
ILCV 3546 → 85\\
ILCV 3598 (em) → 204\\
ILCV 3650 → 126\\
ILCV 3663 (em) → 46\\
ILCV 3696 → 442\\
ILCV 3727a → 346\\
ILCV 3747. C0 conserve 468 textes pour la comparaison historique. C1 inclut la notice 406 : 469 textes, dont 196 avec in pace (41,8 \%). Le statut final et l’événement daté exigent un témoin → 406\\
ILCV 3754a (em) → 209\\
ILCV 3754b (em) → 217\\
ILCV 3781 → 411\\
ILCV 3797a → 24\\
ILCV 3797b (em) → 130\\
ILCV 3800a → 374\\
ILCV 3808 (em) → 141\\
ILCV 3811a → 478\\
ILCV 3813a → 150\\
ILCV 3821 (em) → 340\\
ILCV 3822 → 365\\
ILCV 3904 → 72\\
ILCV 3996c → 206\\
ILCV 4144b → 62\\
ILCV 4144b → 247\\
ILCV 4146 → 194\\
ILCV 4146a → 207\\
ILCV 4146c → 235\\
ILCV 4146d (add, em) → 300\\
ILCV 4181 (em) → 439\\
ILCV 4193 → 414\\
ILCV 4213a → 135\\
ILCV 4215 → 324\\
ILCV 4217a → 114\\
ILCV 4219 → 193\\
ILCV 4219a → 297\\
ILCV 4226 (em) → 184\\
ILCV 4248 → 265\\
ILCV 4248a (em) → 218\\
ILCV 4249a → 286\\
ILCV 4258 → 147\\
ILCV 4258a → 185\\
ILCV 4289 → 200\\
ILCV 4311 → 115\\
ILCV 4326 → 26\\
ILCV 4333a → 237\\
ILCV 4338f → 78\\
ILCV 4351 → 146\\
ILCV 4368 → 60\\
ILCV 4369 (add) → 181\\
ILCV 4379 → 287\\
ILCV 4381 → 349\\
ILCV 4392 → 169\\
ILCV 4393 → 7\\
ILCV 4393a → 128\\
ILCV 4399 (em) → 55\\
ILCV 4399a → 131\\
ILCV 4400a → 416\\
ILCV 4400b → 431\\
ILCV 4411 (em) → 91\\
ILCV 4414 → 359\\
ILCV 4415 (add, em) → 466\\
ILCV 4458 → 168\\
ILCV 4460 → 275\\
ILCV 4461a → 164\\
ILCV 4539 → 428\\
ILCV 4556 → 357\\
ILCV 4606 → 102\\
ILCV 4620 (em) → 262\\
ILCV 4623 (em) → 298\\
ILCV 4642 → 355\\
ILCV 4647 → 1\\
ILCV 4653 → 74\\
ILCV 4670 → 155\\
ILCV 4671 → 399\\
ILCV 4700 → 231\\
ILCV 4703 → 472\\
ILCV 4744 (em) → 470\\
ILCV 4763 → 266\\
ILCV 4768 → 397\\
ILCV 4988 (em) → 409\\
ILMN-01, 440 → 64\\
ILMN-01, 441 → 78\\
ILMN-01, 442 → 79\\
ILMN-01, 443 → 128\\
ILMN-01, 444 → 146\\
ILMN-01, 447 → 405\\
IMCCatania 00360 → 99\\
InscrIt-04-01, 543 → 439\\
Louvre 00852 → 80\\
Louvre 00966 → 5\\
Luciani-02, 18 → 118\\
Mossong-2022, 27 → 321\\
Mossong-2022, 28 → 378\\
Mossong-2022, 30 → 400\\
Mossong-2022, 127 → 105\\
Mossong-2022, 129 → 134\\
Mossong-2022, 131 → 289\\
Mossong-2022, 132 → 310\\
Mossong-2022, 133 → 428\\
Pittori 00008 → 229\\
RAC-1924-82 → 20\\
RAC-1925-14 → 416\\
RAC-1925-24,1 → 72\\
RAC-1981-223 → 72\\
RIRIOe 00095 → 276\\
SEG-30, 1208 → 447\\
Trout-01, 51 → 244\\
Veleia-2012-31 → 244\\
aucun mot transmis n’est modifié → 406\end{multicols}

